% Standalone audit manuscript. All sections, appendices, and bibliography are embedded.
\documentclass[11pt]{article}
\usepackage{iftex}
\ifPDFTeX
  \usepackage[T1]{fontenc}
  \usepackage[utf8]{inputenc}
\fi
\usepackage{amsmath,amssymb,amsthm,mathtools}
\usepackage{mathrsfs}
\usepackage[margin=28mm]{geometry}
\usepackage{booktabs,longtable,array}
\usepackage[hidelinks]{hyperref}
\usepackage{microtype}
\setlength{\emergencystretch}{2em}
\raggedbottom
\hypersetup{pdftitle={A candidate natural-boundary proof for the bulk square-lattice Ising susceptibility},pdfauthor={Dehao Lin},pdfsubject={AI-assisted candidate proof, v0.1-rc4; candidate for public mathematical review; not independently human- or proof-assistant-certified}}
\numberwithin{equation}{section}
\newtheorem{theorem}{Theorem}[section]
\newtheorem{lemma}[theorem]{Lemma}
\newtheorem{proposition}[theorem]{Proposition}
\newtheorem{corollary}[theorem]{Corollary}
\theoremstyle{definition}
\newtheorem{definition}[theorem]{Definition}
\newtheorem{hypothesis}[theorem]{Hypothesis}
\newtheorem{remark}[theorem]{Remark}
\newcommand{\dd}{\,\mathrm d}
\newcommand{\ii}{\mathrm i}
\newcommand{\T}{\mathbb T}
\newcommand{\R}{\mathbb R}
\newcommand{\C}{\mathbb C}
\newcommand{\Q}{\mathbb Q}
\newcommand{\Z}{\mathbb Z}
\newcommand{\eps}{\varepsilon}
\newcommand{\M}{\mathcal M}
\newcommand{\Pf}{\operatorname{Pf}}
\newcommand{\Res}{\operatorname{Res}}
\newcommand{\Li}{\mathcal L}
\newcommand{\dist}{\operatorname{dist}}
\newcommand{\sstar}{s_*}
\newcommand{\se}{s_\varepsilon}
\newcommand{\cO}{\mathcal O}
\newcommand{\gap}[1]{\par\medskip\noindent\textbf{SOURCE-SUPPORTED PROOF GAP:} #1\par\medskip}
\title{A candidate natural-boundary proof\\
for the bulk square-lattice Ising susceptibility}
\author{Dehao Lin\\
\small School of Physics, Sun Yat-sen University\\
\small\texttt{lindh9@mail2.sysu.edu.cn}}
\date{Version 0.1-rc4 (candidate for mathematical review)\\1 October 2026}
\begin{document}
\maketitle

% BEGIN SOURCE: sections/01_introduction.tex
\begin{abstract}
This manuscript presents a candidate proof that the unit circle is a
natural boundary for the zero-field, isotropic, infinite square-lattice
Ising susceptibility continued from low temperature. The argument has
two separate analytic tasks. At a dense cyclotomic family of boundary
points, an actual average-angle residue produces a nonzero first
singular derivative. A weighted contour deformation and a partition of
the entire particle-number tail are then used to claim an absolute
smaller-order bound for every derivative needed in the bulk product
rule. The decisive estimates, support conditions, differentiation
recurrences, and order of choices are included in the paper.

This release candidate is reconstructed from an AI-generated research checkpoint.
Its propositions are the proof claims to be checked, not independently
certified results. Appendix~\ref{app:auditmap} identifies the load-bearing
checks and the coverage of all tail sectors. No exhaustive priority
clearance is claimed; the natural-boundary conclusion has no
independent certification.
\end{abstract}

\section{The claim, scope, and proof architecture}\label{sec:scope}
Let \(s=\sinh(2\beta_{\mathrm{phys}} J)\), with inverse temperature
\(\beta_{\mathrm{phys}}\), zero magnetic field, and equal positive
nearest-neighbor coupling \(J\) on the infinite square lattice. Write
\[
 \mathcal X(s)=\beta_{\mathrm{phys}}^{-1}\chi(s),\qquad
 \M^2(s)=(1-s^{-4})^{1/4},
\]
The factor \(\M^2\) is the square of the low-temperature spontaneous
magnetization \cite{Yang1952}; its branch tends to one at infinity.
The symbol \(\beta_{\mathrm{phys}}\) distinguishes inverse temperature
from the geometric angle \(\beta\) used below. We study the exterior
germ inherited from real \(s>1\). On that real interval the response is
defined by connected correlations in Palmer--Tracy's infinite-volume
\(+\) state, obtained as the limit of \(+\) boundary conditions
\cite[pp.~332, 361]{PalmerTracy1981}. Global spin flip gives the same
two-point function and the same connected correlation after subtracting
\(\M^2\) in the \(-\) state. The function \(\mathcal X\) is the full
bulk response. The proposed conclusion is
\[
 \text{no point of }|s|=1\text{ admits a local holomorphic
 continuation of this exterior germ.}
\]
The physical factor \(\beta_{\mathrm{phys}}(s)=\operatorname{arsinh}(s)/(2J)\) is locally
analytic and nonzero at each selected nonaxis point on the continuation
under consideration. Its logarithmic monodromy is not replaced by an
assertion of single-valuedness on the entire exterior.

An exterior holomorphic function has \(\T\) as a natural boundary if no
point of \(\T\) has a neighborhood to which its exterior germ extends
holomorphically. A dense set of points at which a finite-order derivative
is unbounded on an exterior radial approach is sufficient.

\subsection{Related work and the proposed contribution}
The form-factor expansion of the bulk susceptibility goes back to
Wu, McCoy, Tracy, and Barouch \cite{WMTB1976}. Later correlation
representations include the low-temperature Fredholm formula of
Palmer and Tracy \cite{PalmerTracy1981}, used in Appendix 1 of
Tracy and Widom \cite{TW2014}. Guttmann and Enting
\cite{GuttmannEnting1996} supplied early numerical evidence for a
natural boundary. Nickel \cite{Nickel1999,Nickel2000} identified
the unit-circle form-factor singularities now bearing his name;
Boukraa et al.\ \cite{Boukraa2008} developed further series and
singularity evidence. These findings motivate the full-sum question.

Orrick, Nickel, Guttmann, and Perk \cite[Section~3.2]{Orrick2001}
examined the order dependence of singularity amplitudes and argued
that it provides strong evidence against cancellation in the full sum.
Their introduction explicitly distinguishes this evidence from a proof.
We cite that analysis as a heuristic precedent, not as an infinite-tail
theorem.

For the diagonal susceptibility, Tracy and Widom \cite{TW2013}
proved a natural-boundary theorem. Their Section IV isolates a
divergent derivative of one particle term and bounds the corresponding
derivative of the remaining sum. Their later work
\cite{TW2017Toeplitz} extends that direction to more general sums
of Toeplitz determinants. Thus the architecture of isolating a
singular derivative and controlling the tail to exclude cancellation
is an established method, not a new principle claimed here.
The analysis in Sections III--IV of \cite{TW2017Toeplitz} also uses
contour partitions, Vandermonde homogeneity, and Hadamard estimates
with factorial suppression; grouping near different singularities
plays a role there as well. These tools and organizing ideas are
precedents for this work, not inventions attributed to the present
manuscript. Their bulk implementation must still be justified here.

The diagonal result concerns correlations summed along a single
lattice diagonal. The present observable sums over the full lattice;
its representation \eqref{eq:double} contains the two global factors
\((1-X)(1-Y)\). The diagonal theorem supplies neither the required
bulk first-chart analysis nor the bulk uniform tail estimate.
Tracy and Widom's separate fixed-order result \cite{TW2014}
gives smooth extension of each even form factor away from its Nickel
set, without excluding cancellation in the infinite bulk sum.

The proposed technical work here is the selected uniquely-first point
family, the compatible whole-integral first-amplitude calculation,
and the estimates for the weighted contour decomposition, protected
parameter disks, and differentiated tail sectors of the bulk integral.
These are the claims to assess against the displayed derivations and
the prior literature. Neither an internal audit verdict nor the
inclusion of additional references establishes their correctness or
their novelty. No exhaustive literature-priority clearance is claimed.

\subsection{The two estimates that carry the conclusion}
For each selected point \(s_*\), Section~\ref{sec:arithmetic} specifies
a first even order \(N_0=2p\) and the fixed integer
\(k=N_0^2/2-1\). Put \(s_\eps=(1+\eps)s_*\). The proposed argument
establishes
\begin{align}
 T_{N_0}^{(k)}(s_\eps)
   &=2L_\beta\eps^{-1/2}+o(\eps^{-1/2}),\qquad L_\beta\ne0,
       \label{eq:roadmapfirst}\\
 \sum_{\substack{N>N_0\\N\ {\rm even}}}
       |T_N^{(j)}(s_\eps)|
   &=o(\eps^{-1/2}),\qquad 0\le j\le k .
       \label{eq:roadmaptail}
\end{align}
The finitely many lower form factors and the lower derivatives of the
first singular form factor remain bounded. The normalization in
Section~\ref{sec:setup} then gives
\[
 \mathcal X^{(k)}(s_\eps)
   =4\M^2(s_*)L_\beta\eps^{-1/2}+o(\eps^{-1/2}).
\]
Section~\ref{sec:conclusion} proves this implication separately.
In particular, a proof only of \eqref{eq:roadmapfirst}, or a tail
estimate only for \(j=k\), does not complete the proposed argument.

\subsection{How the proof is organized}
The arithmetic and first-amplitude arguments are in
Sections~\ref{sec:arithmetic} and \ref{sec:first}. The tail begins
with the exact identity \(T_N=F_N+K_N+S_N\) in
Section~\ref{sec:tailreduction}; \(S_N\) is the actual weighted
Stokes correction. Sections~\ref{sec:F}--\ref{sec:sectors} estimate
every term in that identity. Appendix~\ref{app:lie} supplies the
fixed-domain, high-order differentiation rules used in those estimates.
The constants and quantifiers are collected in
Appendix~\ref{app:constants}; Appendix~\ref{app:auditmap} is a route
through the proof for an adversarial reader.

Throughout, the point and the differentiation order are fixed before
\(\eps\) tends to zero. Constants may depend on both. Uniformity in
the primes indexing the dense family is neither asserted nor needed.
A bound uniform in \(N\) always refers to the explicitly stated
\(N\)-window. The contour radius is held fixed during differentiation;
window boundaries depending on \(\eps\) are introduced only afterwards.

\begin{center}
\begin{tabular}{p{.20\textwidth}p{.71\textwidth}}
\toprule
Symbol & Meaning\\
\midrule
\(N,N_0,k\) & Even particle number; first singular order; fixed derivative order.\\
\(S,W,z,\phi\) & Dispersion variables, with \(z=e^{-\ii\phi}\) on the lower sheet.\\
\(Y,Z,P_{ij}\) & Products of the \(y,z\) coordinates and the complete canceled pair.\\
\(p,m,a,P\) & Fixed deformation/selector functions and occupancy \(P=\sum_i p_i\).\\
\(H,\lambda_*\) & \(\log(1/\eps)\) and the post-differentiation current split \(\eps^{\alpha_0}\).\\
\(b_N,\rho_N\) & Exponential microcore and full-equality scales in the intermediate window.\\
\bottomrule
\end{tabular}
\end{center}

% END SOURCE

\tableofcontents

% BEGIN SOURCE: sections/02_setup.tex
\section{The normalized bulk representation}\label{sec:setup}
% PROV-INTERNAL: FORM_FACTOR_NORMALIZATION_AUDIT.md; ONSITE_GROUPING_EXACT_AUDIT.md.
% PROV-PRIMARY: TW2014 Appendix 1 (7), normalized auxiliary residue and lattice sum.
% RECONSTRUCTION: retain every normalized integration measure.
Put \(S(s)=s+s^{-1}\) and
\[
 D(x,y;s)=S(s)-\tfrac12(x+x^{-1})-\tfrac12(y+y^{-1}).
\]
We take as external physical input the low-temperature correlation
expansion in the published Appendix~1, equations~(7)--(8),
pp.~1132--1133 of \cite{TW2014}. Its upstream source is
\cite[Cor.~2.0, p.~351; Thm.~4.1, p.~361; Thm.~5.0, p.~374,
and proof pp.~375--378]{PalmerTracy1981}.
For even \(N\), let \(T_N\) denote the Fredholm-normalized offsite
form factor. With positively oriented circles sufficiently close to the
unit circle, it is
\begin{align}
 T_N(s)&=\frac1{N!}\oint\cdots\oint
 \frac{X^{-1}+Y^{-1}}{(1-X)(1-Y)}
 \prod_{i<j}\frac{x_i-x_j}{1-x_ix_j}
              \frac{y_i-y_j}{1-y_iy_j}
 \prod_{i=1}^N\frac{\dd x_i\,\dd y_i}{(2\pi\ii)^2D(x_i,y_i;s)},\label{eq:double}\\
 X&=\prod_i x_i,\qquad Y=\prod_i y_i, \nonumber\\
 \mathcal X(s)&=1-\M^2(s)+2\M^2(s)\sum_{\substack{N\ge2\\N\ {\rm even}}}T_N(s).
 \label{eq:bulk}
\end{align}
This normalization follows by converting the \(N\) normalized angular
measures in Appendix~1, equation~(7), of \cite{TW2014} and retaining
the \(N\) normalized auxiliary residues. It agrees with the prefactor
\((2\pi\ii)^{-2N}/N!\) displayed with the published equation~(3).
That coefficient sums all sites, including the origin, and equals
\(f_{00}^{(N)}+2T_N\) at fixed even \(N\).
Appendix~\ref{app:normalization} defines these terms and derives the relation.

% PROV-INTERNAL: RESIDUE_REDUCTION.md, lines 9-42.
% PROV-PRIMARY: TW2014 (2), Appendix 1 (7) and the normalized residue identity.
% RECONSTRUCTION: successive residues and subtraction of dispersion equations.
\begin{lemma}[Residue reduction]\label{lem:residue}
Let
\[
 W(y;s)=S(s)-\tfrac12(y+y^{-1}),\quad
 z=W-\sqrt{W^2-1}=e^{-\gamma},\quad
 R(y;s)=\frac{2z^2}{1-z^2}.
\]
Choose the square root by continuation of the interior root and require
\(\lvert z(y;s)\rvert<r<1\) on the \(y\)-circle. Then
\begin{equation}
 T_N=\frac1{N!}\oint\cdots\oint
 \frac{Z^{-1}+Y^{-1}}{(1-Z)(1-Y)}
 \prod_{i<j}P_{ij}\prod_i R(y_i;s)\frac{\dd y_i}{2\pi\ii},
 \qquad Z=\prod_i z_i,\label{eq:reduced}
\end{equation}
where
\begin{equation}
 P_{ij}=\frac{z_i-z_j}{1-z_iz_j}\frac{y_i-y_j}{1-y_iy_j}
       =-\frac{(y_i-y_j)^2z_iz_j}{y_iy_j(1-z_iz_j)^2}.
 \label{eq:pair}
\end{equation}
\end{lemma}
\begin{proof}
The only enclosed dispersion root in each \(x_i\) is \(z_i\);
its residue is \(2z_i^2/(1-z_i^2)\). At \(x_i=0\) the reciprocal dispersion
factor is \(-2x_i+O(x_i^2)\), canceling the apparent numerator pole.
All pair and product poles involving the other interior \(x\)'s are
outside the \(x_i\)-circle. These statements remain true after each
successive residue because each previously substituted root has modulus
less than \(r\).
Subtracting \(z_i+z_i^{-1}=2S-y_i-y_i^{-1}\) for two indices gives
\[
 (z_i-z_j)(1-z_iz_j)
 =-\frac{z_iz_j}{y_iy_j}(y_i-y_j)(1-y_iy_j),
\]
which proves the canceled expression. Each residue consumes its own
normalized \(x\)-measure and leaves a normalized \(y\)-measure.
\end{proof}

For a fixed \(\sstar=e^{\ii\theta_*}\), \(0<\theta_*<\pi/2\), set
\(\se=(1+\eps)\sstar\). On a disk \(\lvert s-\se\rvert<c_1\eps\)
choose \(r=e^{-c_0\eps}\), with \(c_0,c_1>0\) sufficiently small.
Then \(\Im W\ge c_2\eps\), uniformly in the angle. The inverse root is
in the lower unit half-disk. Writing \(\gamma=u+\ii v\) gives
\(\Im W=\sinh u\sin v\le\sinh u\), hence \(u\ge c_3\eps\).
Choosing \(c_0<c_3\) verifies the stronger condition \(|z|<r\), not merely
\(|z|<1\). During an \(s\)-derivative this chosen \(r\) is held fixed.

% PROV-INTERNAL: TWO_NONZERO_PRIME_FAMILY.md, lines 46-50;
% FORM_FACTOR_SIGN_SYMMETRY.md.
% PROV-PRIMARY: TW2014 contour representation and Appendix 1.
\begin{lemma}[Symmetries]\label{lem:symmetry}
The exterior normalized susceptibility and its even form factors are even
in \(s\) and invariant under conjugation of the argument and value.
\end{lemma}
\begin{proof}
In \eqref{eq:double} replace \(s,x_i,y_i\) by \(-s,-x_i,-y_i\).
Each dispersion factor changes sign; their number is even. Each pair
changes sign in both variables and each product is unchanged.
The contour substitution preserves the integral. The branch of \(\M^2\)
is also even. Conjugation follows either by conjugating these contours
with their orientations and measures together or from the real
low-temperature Fredholm expansion and the identity theorem.
\end{proof}

% END SOURCE


% BEGIN SOURCE: sections/03_dense_nickel_family.tex
\section{A dense uniquely-first family}
\label{sec:arithmetic}
For this paper a unit-circle point is a Nickel point of order \(n\)
when
\[
 s_*+s_*^{-1}
   =\cos(2\pi l/n)+\cos(2\pi m/n)
 \quad\text{for some integers }l,m.
\]
This is Nickel's singularity condition
\cite{Nickel1999,Nickel2000}, in the convention of \cite{TW2014}.
Only even orders occur in the low-temperature expansion under study.
The theorem below establishes the additional first-order and uniqueness
properties for the specific point family used here.
% PROV-INTERNAL: TWO_NONZERO_PRIME_FAMILY.md, lines 5-50.
% PROV-PRIMARY: TW2014 definition of a Nickel point.
% RECONSTRUCTION: elementary cyclotomic argument; no Durham notes used.
\begin{theorem}\label{thm:prime}
Let \(p\ge11\) be prime, and choose distinct integers \(a,b\) with
\(\alpha=2\pi a/p,\ \beta=2\pi b/p\in(0,\pi/2)\).
Set
\[
 c=\frac{\cos\alpha+\cos\beta}{2},\qquad
 \sstar=e^{\ii\arccos c},\qquad N_0=2p.
\]
The first even Nickel order at \(\sstar\) is \(N_0\).
At that order the unordered cosine pair is uniquely
\(\{\cos\alpha,\cos\beta\}\).
These points are dense in the open upper-right arc.
If \(\theta_B=\arccos(2c-1)\), then \(e^{-\ii\theta_B}\) is algebraic
and is not a root of unity.
\end{theorem}
\begin{proof}
Let \(\zeta=e^{2\pi\ii/p}\), \(d=(p-1)/2\),
\(u_j=\zeta^j+\zeta^{-j}\), and \(E=\Q(\zeta+\zeta^{-1})\).
The trace of every nonzero-index \(u_j\) is \(-1\), by summing over the
pairs of nonzero residue classes. Thus
\(\operatorname{Tr}_E(2c)=-1\). If \(2c\) were rational it would be
\(-1/d\), whereas it is positive.

If an even \(n<2p\) represented the same Nickel point, then \(p\nmid n\)
and \(c\in\Q(\zeta_p)\cap\Q(\zeta_n)=\Q\);
the last equality is proved in Lemma~\ref{lem:cyclotomic}.
This contradicts the preceding trace. Order \(2p\) contains the indicated
pair of angles and is therefore first.

At order \(2p\) every cosine is \(\pm u_j/2\) or \(\pm1\).
The only rational relation between \(1,u_1,\ldots,u_d\) is a multiple of
\(1+\sum_j u_j=0\): lift a relation to a polynomial of degree at most
\(p-1\) vanishing at \(\zeta\), and use its minimal polynomial
\(1+X+\cdots+X^{p-1}\). An equality between two cosine pairs uses at most
four of the \(u_j\). Since \(d\ge5\), one coefficient is zero, forcing
that multiple to vanish. The remaining coefficients give precisely the
two positive terms \(u_a,u_b\), with no constant term.

For any \(c_0\in(0,1)\), along unbounded primes choose two distinct
neighboring grid angles tending to \(\arccos c_0\). They eventually lie
in \((0,\pi/2)\), and their cosine averages tend to \(c_0\).

Finally \(w=2\cos\theta_B=u_a+u_b-2\) has trace \(-2d-2\).
Were \(e^{\ii\theta_B}\) a root of unity, every conjugate of \(w\)
would be in \([-2,2]\); every embedding of \(E\) restricts to such a
conjugate. Its trace could not be less than \(-2d\).
Algebraicity follows from \(x^2-wx+1=0\).
\end{proof}

% PROV-INTERNAL: NONRESONANT_BRANCH_EXPONENTIAL_CUTOFF.md, lines 5-21.
% RECONSTRUCTION: nonzero integer resultant and a product estimate.
\begin{lemma}[Exponential separation]\label{lem:resultant}
For the fixed algebraic number \(\xi=e^{-\ii\theta_B}\) there is
\(A>0\) such that \(\lvert1-\xi^N\rvert\ge e^{-AN}\) for every \(N\ge1\).
\end{lemma}
\begin{proof}
Write its primitive integer minimal polynomial as
\(P(X)=a_d\prod_{\nu=1}^d(X-\xi_\nu)\), with \(\xi_1=\xi\).
It has no root in common with \(X^N-1\): such a conjugate root would make
its irreducible polynomial cyclotomic and would make \(\xi\) a root of
unity. Consequently
\[
 1\le |\Res(P,X^N-1)|
 =|a_d|^N\prod_\nu|\xi_\nu^N-1|
 \le |a_d|^N|1-\xi^N|\,2^{d-1}
       \prod_{\nu>1}\max(1,|\xi_\nu|)^N.
\]
Increasing \(A\) absorbs all fixed factors for \(N\ge1\).
\end{proof}

% END SOURCE


% BEGIN SOURCE: sections/04_first_singular_amplitude.tex
\section{The first singular form factor}
\label{sec:first}
Fix the family of Theorem~\ref{thm:prime} and put
\[
 N=N_0,\qquad k=N^2/2-1,\qquad
 S_*=\cos\alpha+\cos\beta=2\cos\theta_*.
\]
All constants in this section are allowed to depend on this fixed \(N\).
At the ordered lower chart \(y_0=e^{-\ii\alpha}\),
\(z_0=e^{-\ii\beta}\), write \(y_j=r e^{\ii(-\alpha+u_j)}\) and
\(z_j=e^{-\ii\phi_j}\). Define
\[
 v=\sum_j u_j,\qquad t_j=u_j-v/N,\qquad \sum_jt_j=0.
\]
The absolute Jacobian of \((v,t_1,\ldots,t_{N-1})\mapsto(u_1,\ldots,u_N)\)
is one.

\paragraph{Compatible localization.}
Let \(\Omega^{\mathrm{dbl}}_{N,s}\) denote the complete normalized
double-contour density in \eqref{eq:double}. Around each of the two
all-equal lower \(y\)-configurations, choose a smooth angular weight
\[
 w_\gamma(\theta)=\eta_\gamma(v_\gamma)\chi_\gamma(t_\gamma),
 \qquad \gamma\in\{\alpha,\beta\}.
\]
Here \(u_{\gamma,j}=\theta_j+\gamma\),
\(v_\gamma=\sum_j u_{\gamma,j}\), and
\(t_{\gamma,j}=u_{\gamma,j}-v_\gamma/N\).
Choose \(\eta_\gamma=1\) on \(|v_\gamma|\le\delta/2\), supported
in \(|v_\gamma|<\delta\), and \(\chi_\gamma=1\) near the zero shape,
supported in a fixed small shape ball. Extend each weight smoothly by
zero on the angular torus, with disjoint supports and values in \([0,1]\).
The weights depend only on the \(y\) angles and are independent of
\(s\), the \(x\) angles, and the varying evaluation radius.
Thus the exact identity
\[
 T_N=I_\alpha+I_\beta+I_{\mathrm{rem}},\qquad
 I_\gamma=\int w_\gamma\Omega^{\mathrm{dbl}}_{N,s},\qquad
 I_{\mathrm{rem}}=\int(1-w_\alpha-w_\beta)
                         \Omega^{\mathrm{dbl}}_{N,s}
\]
holds before any localization in \(x\).
The remaining weight vanishes on an open neighborhood of both bad
double-torus configurations. The classification and estimate in
Lemma~\ref{lem:complement} therefore give
\(\partial_s^j I_{\mathrm{rem}}=O(1)\) for every fixed \(j\).
Its further smooth partition may depend on both sets of angles,
since no residue calculation is applied to these complementary pieces.

In each \(I_\gamma\), all \(x\) integrations remain unweighted.
Lemma~\ref{lem:residue} consequently gives the reduced integral with
the same \(y\)-weight. In this reduced local integral replace
\(\eta_\gamma(v)\) by the indicator of \([-\delta,\delta]\).
Their difference is supported where \(\delta/2\le |v|\le\delta\).
Using the coefficients \(b,d\) computed below, choose the shape radius
\(h\) so that \(dh^2\ll b\delta\), and then take \(\delta\) and
the parameter neighborhood sufficiently small. On this edge region,
the \(Y\) phase is separated by \(|v|\), while the distance of
\(\Re\sum_j\phi_j\) from \(2\pi\Z\) is bounded below by
\(b|v|-C(h^2+v^2+\eps)\). Both global denominators, and every local
one-body and pair denominator, are therefore uniformly separated
from zero. The replacement changes each fixed \(s\)-derivative by
\(O(1)\), with the radius held fixed during differentiation.
Only this hard mean segment is moved in Lemma~\ref{lem:mean}; the
shape cutoff stays fixed. No nonholomorphic angular weight is moved
through an \(x\) residue or a complex mean deformation.

% PROV-INTERNAL: COMPACT_PRIME_MEAN_RESIDUE_AMPLITUDE.md, lines 5-41, 55-69.
% PROV-PRIMARY: TW2014 Appendix 1; no printed amplitude is imported.
\begin{lemma}[Physical mean residue]\label{lem:mean}
Choose \(r=e^{-c_0\eps}\) as in Lemma~\ref{lem:residue}. The singular
part of the localized chart is obtained by moving the real \(v\)-segment
downward through \(Y=1\); the resulting coefficient is \(+2\pi\).
At that residue the denominator satisfies
\[
 |1-Z(t,\se)|\ge c_1\bigl(\eps+\textstyle\sum_jt_j^2\bigr)
\]
on a sufficiently small fixed shape ball.
\end{lemma}
\begin{proof}
Implicit differentiation of \(\cos\phi=S-\cos(-\alpha+u)\) gives
\[
 b=\phi_u=\frac{\sin\alpha}{\sin\beta}>0,\qquad
 d=-\tfrac12\phi_{uu}
   =\frac{S_*(1-\cos\alpha\cos\beta)}{2\sin^3\beta}>0.
\]
At the radial point,
\[
 \phi=\beta-\ii A_0\eps+bu-du^2
       +O(\eps^2+\eps|u|+|u|^3),\quad
 A_0=\frac{2\sin\theta_*-c_0\sin\alpha}{\sin\beta}.
\]
Choose \(c_0\) so that \(A_0>c_0\), in addition to global residue
admissibility. On a rectangle \(|\Re v|\le\delta,\ -h_0\le\Im v\le0\),
take the fixed shape support \(\sum t_j^2<h^2\) with \(dh^2\ll b\delta\).
Take \(h_0\) sufficiently small but fixed and \(Nc_0\eps<h_0\).
The \(Y\)-pole is \(v=-\ii Nc_0\eps\). The \(Z\)-pole is above the
real mean line in the local model; throughout this rectangle a downward
mean motion increases its damping. The side phases are separated by
\(b\delta+O(h^2+h_0^2)\), and the bottom has fixed positive damping.
All pair denominators remain nonzero since \(z_0^2\ne1\).
Thus the three displaced sides, with any fixed number of \(s\)-derivatives,
are bounded. A shape cutoff independent of \(v\) is legitimate here;
smoothing the hard mean endpoints is supported on the separated sides.

The downward closure is clockwise. Since
\(\partial_v(1-Y)=-\ii\) at the pole, the orientation and derivative
give \((-2\pi\ii)/(-\ii)=+2\pi\).
There each \(y_j=e^{\ii(-\alpha+t_j)}\) lies exactly on the unit circle.
On this compact regular arc, \(-\log|Z|\ge cN\eps\).
The real phase \(g_\eps(\theta)=\Re\arccos(S(\se)-\cos\theta)\) has
\(g_\eps''\le-c<0\), and
\(N g_\eps(-\alpha)=N\beta+O(\eps^2)\).
Taylor's integral remainder on \(\sum t_j=0\) supplies a phase
separation \(c\sum t_j^2\) when that sum exceeds a sufficiently large
multiple of \(\eps^2\). Below this scale the modulus separation
controls it. Shrinking the ball prevents phase wrapping. Combining
the two separations proves the asserted nonlinear bound.
\end{proof}

% PROV-INTERNAL: COMPACT_PRIME_MEAN_RESIDUE_AMPLITUDE.md, lines 25-53.
% RECONSTRUCTION: residue first, then an absolutely convergent beta integral.
\begin{lemma}[Nonzero shape period]\label{lem:period}
Let
\[
 a_0=\frac{N^2-1}{2},\quad Q=\frac{2N\sin\theta_*}{\sin\beta},\quad
 \mathscr C_N=\frac1{\sqrt N}
 \int_{\{\omega\in\R^N:\,\sum\omega_j=0,\ |\omega|=1\}}
       \Delta(\omega)^2\,\dd\sigma(\omega)>0 .
\]
Then
\begin{equation}
 J_\beta=\int_{\R^{N-1}}\frac{\Delta(\tau)^2\,
                  \dd\tau_1\cdots\dd\tau_{N-1}}
                   {(Q-\ii d\sum_j\tau_j^2)^{k+1}}
 =\frac{\mathscr C_N}{2}Q^{-1/2}(-\ii d)^{-a_0}
                     B(a_0,1/2)\ne0. \label{eq:beta}
\end{equation}
Here \(\tau_N=-\sum_{j<N}\tau_j\), and the power is obtained by
continuation from a quadratic coefficient with positive real part.
\end{lemma}
\begin{proof}
Polar coordinates in the zero-sum hyperplane give radial numerator
\(R^{N^2-2}\); the Euclidean-to-coordinate volume factor is \(1/\sqrt N\).
The substitution \(u=R^2\) gives Euler's beta integral.
Continuation to the coefficient \(-\ii d\) is dominated by a constant
times \(R^{N^2-2}(1+R^2)^{-N^2/2}\), whose tail is \(O(R^{-2})\).
All factors on the right of \eqref{eq:beta} are nonzero.
\end{proof}

The leading local density and pole derivative are
\[
 K_\beta=\frac2{N!}\,2^{-N(N-1)}(2\pi)^{-N}\sin(\beta)^{-N^2},
 \qquad
 D_s=-\frac{\ii N S'(\sstar)}{\sin\beta}
     =\frac{2N e^{-\ii\theta_*}\sin\theta_*}{\sin\beta}.
\]
Indeed every pair has leading factor
\(- (u_i-u_j)^2/(4\sin^2\beta)\), and each residue-angle measure
has leading factor \(y_0z_0/(2\pi\ii\sin\beta)\).
Since \(N\) is even and \(y_0^Nz_0^N=1\), their combined phase is
\(\ii^{-N^2}=1\), giving \(K_\beta>0\).

After the actual mean residue, put \(t=\sqrt\eps\,\tau\).
The \(k\)-th derivative has leading coefficient
\[
 L_\beta=2\pi K_\beta(-1)^k k!D_s^kJ_\beta\ne0
\]
times \(\eps^{-1/2}\). The bound in Lemma~\ref{lem:mean} dominates
the rescaled integrand by
\(C\Delta(\tau)^2/(1+|\tau|^2)^{k+1}\).
Lower pole orders vanish after multiplication by \(\sqrt\eps\).
This is a dominated limit of the constrained shape integral,
not of the joint unintegrated mean/shape absolute value.
The sine powers in \(K_\beta D_s^kQ^{-1/2}d^{-a_0}\) sum to
\(-N^2-k+1/2+3a_0=0\). Hence \(L_\alpha=L_\beta\)
for the swapped ordered chart, including the phase. Exchanging all
\(x\)'s with all \(y\)'s in \eqref{eq:double} also preserves orientation
for even \(N\) and checks this equality before asymptotics.

% PROV-INTERNAL: COMPACT_PRIME_MEAN_RESIDUE_AMPLITUDE.md, lines 73-79.
% PROV-PRIMARY: TW2014 Section 3, proof of Lemmas 1-3 and Appendix 2.
\begin{lemma}[Complement of the two first charts]\label{lem:complement}
At \(N=N_0\) all other localized charts of the original double contour
have bounded derivatives of every fixed order on the exterior radial
approach.
\end{lemma}
\begin{proof}
The local singular-factor gradients in \cite{TW2014} are the two
all-ones product vectors, pair vectors, and the dispersion vectors with
entries \((\Im x_j,\Im y_j)\).
A nonnegative zero combination cannot contain a pair vector:
cancellation of its two positive entries would require two negative
imaginary parts, whereas reciprocal coordinates have opposite imaginary
parts. If both product vectors occur, every \(x_j,y_j\) is lower,
the ratios \(\Im x_j/\Im y_j\) are equal, and the dispersion relation
forces all \(x_j\)'s equal and all \(y_j\)'s equal. The strictness follows
by differentiating that ratio along
\(\cos\theta+\cos\varphi=S_*>0\). Both common values are \(N\)-th
roots of unity. Theorem~\ref{thm:prime} leaves exactly the two charts.
If only one product vector occurs, the other coordinates are all \(+1\)
or all \(-1\). The former would make \(e^{-\ii\theta_B}\) an \(N\)-th
root of unity; the latter requires a real cosine \(S_*+1>1\).
Both are excluded.

For a fixed smooth weight whose support avoids these two
configurations, cover its support by finitely many compact localized
pieces. On each piece choose a real separating direction \(e\) and
\(\delta_*>0\) with \(e\cdot v_i\ge\delta_*\) for every active
singular-factor vector \(v_i\). Shrink the piece once so that this
separation persists uniformly for the nearby damped contours.

Represent its \(\ell\) potentially singular factors by damped
positive-half-line exponential integrals, and put
\(R=\sum_{i=1}^{\ell}\xi_i\) for their nonnegative auxiliary
parameters. The exponent \(g\) satisfies \(\Re g\le0\),
\(|e\cdot\nabla g|\ge cR\), and every fixed angular derivative
of \(g\) is \(O(R)\). An \(s\)-derivative of order \(j\), taken
at fixed contour radius, adds a polynomial in the auxiliary parameters
of degree at most \(j\), with uniformly bounded smooth coefficients.
On \(R\ge1\), integrate by parts \(q=\ell+j+1\) times using
\[
 \mathcal L=\frac{e\cdot\nabla}{e\cdot\nabla g},
 \qquad \mathcal L e^g=e^g.
\]
The transpose of this operator gains one power \(R^{-1}\) per
iteration, including derivatives of its coefficients and the localized
weight. The angular integral is therefore bounded by \(CR^{j-q}\).
The auxiliary radial volume is \(R^{\ell-1}\dd R\), so the
remaining tail is bounded by \(C\int_1^\infty R^{-2}\dd R\).
The region \(R\le1\), and pieces with no singular factor, are bounded
directly. The finite cover supplies common constants for each fixed
\(N,j\). This proves in particular the bound for the remaining
weight in the compatible localization above. It uses the local proof
mechanism, not the non-Nickel hypothesis of the global theorem at
the selected Nickel point.
\end{proof}

% PROV-INTERNAL: COMPACT_PRIME_MEAN_RESIDUE_AMPLITUDE.md.
% RECONSTRUCTION: Lemmas above; no extrapolation of a higher-index coarea bound.
\begin{theorem}\label{thm:first}
At the selected point,
\[
 T_{N_0}^{(k)}(\se)=2L_\beta\,\eps^{-1/2}+o(\eps^{-1/2}),
 \qquad L_\beta\ne0.
\]
Every derivative of order \(0\le j<k\) of this first term is bounded
on the radial approach. Every even \(N<N_0\) has bounded derivatives
there as well.
\end{theorem}
\begin{proof}
The compatible \(y\)-only partition gives the two reduced local
integrals and a complement bounded by Lemma~\ref{lem:complement}.
Replacing their smooth mean cutoffs by hard segments adds only bounded
edge terms. The two actual mean residues therefore add with the
displayed coefficients. For lower derivatives the constrained radial
density at \(\eps=0\)
has exponent \(N^2-2j-4\ge0\), so it is integrable.
The displaced sides are bounded as before. Lower even orders are
non-Nickel by Theorem~\ref{thm:prime}, and the fixed-order smoothness
theorem of \cite{TW2014} applies.
\end{proof}

% END SOURCE


% BEGIN SOURCE: sections/05_uniform_tail_reduction.tex
\section{Exact contour decomposition and quantitative requirements}
\label{sec:tailreduction}
For the remainder fix \((p,a,b)\), \(N_0=2p\), and \(k=N_0^2/2-1\).
Set \(H=\log(1/\eps)\) and choose
\begin{equation}
 0<\alpha_0<\frac1{2(k+2)},\qquad \lambda_*=\eps^{\alpha_0}.
 \label{eq:alphachoice}
\end{equation}
Neither \(N_0\) nor \(k\) varies with \(\eps\).
Let \(\Omega_{N,s}\) be the reduced holomorphic form in
\eqref{eq:reduced}, including its \(1/N!\).
Choose smooth periodic functions \(p,m,a\) with values in \([0,1]\)
as follows, in angular representatives \(-\pi<\theta\le\pi\).
The closed support of \(m\) is contained in a small neighborhood of
\(-\theta_B\) within \((-\pi,0)\); \(m=1\) on an open neighborhood
of the closed branch-cutoff supports in Lemma~\ref{lem:partition}.
Thus it is one on the supports of all their finite jets as well.
Choose a small fixed endpoint width \(\delta>0\) and take
\[
 a=0\ \hbox{off }[\delta,\pi-\delta],\qquad
 a=1\ \hbox{on }[2\delta,\pi-2\delta].
\]
Make \(\delta\) small enough that the reflection of
\(\operatorname{supp}m\) lies in the interior of the latter interval.
Take \(p\) with closed support in \((0,\pi)\), equal to one on
an open neighborhood of \([\delta,\pi-\delta]\).
Smooth step functions give all these choices simultaneously.
Their widths and the deformation strength are fixed in the order
of Appendix~\ref{app:constants}, before \(\eps\to0\).
The supports of \(p,m\) are disjoint, the upper inverse branch is
inside \(a=p=1\), and \(\operatorname{supp}a\) and
\(\operatorname{supp}a'\) have positive distance from both endpoints.
In particular a named upper angle has \(\sin\theta\ge\sin\delta>0\).

Define, with fixed deformation strength \(\tau>0\),
\[
 P=\sum_i p(\theta_i),\quad
 h_i=-2\tau p(\theta_i)+\frac{\tau P}{2N}m(\theta_i),\quad
 y_i(\lambda,\theta)=r e^{\ii\theta_i+\lambda h_i},
 \qquad \psi(\theta)=\prod_i(1-a(\theta_i)).
\]
The letter \(P\) here is an occupancy sum, not the pair factor \(P_{ij}\).
Write \(U_{\lambda,s}=\Phi_\lambda^*\Omega_{N,s}\).

% PROV-INTERNAL: GLOBAL_DEFORMATION.md; GENERIC_UPPER_ANCHOR_AUDIT.md;
% UPPER_CUTOFF_STOKES_RETRACTION.md; STOKES_TOP_FORM_ORIENTATION_REDERIVATION.md.
% RECONSTRUCTION: exact holomorphic homotopy, matrix determinant, Cartan identity.
\begin{proposition}[Exact weighted retraction]\label{prop:selector}
For exterior \(s\) in its admissible neighborhood,
\begin{align}
 T_N&=F_N+K_N+S_N,\label{eq:FKS}\\
 F_N&=\int(1-\psi)U_{1,s},\qquad K_N=\int\psi U_{0,s},\nonumber\\
 S_N&=-2\ii\tau\int_0^1\sum_q\int
                         (\partial_{\theta_q}\psi)U_{\lambda,s}\dd\lambda.
 \label{eq:current}
\end{align}
The current has one named index \(q\) in \(\operatorname{supp}a'\),
with \(p_q=1,\ p'_q=m_q=m'_q=0\); all other indices belong to
the enlarged lower support \(K=\operatorname{supp}(1-a)\).
\end{proposition}
\begin{proof}
We have \(\sum_i h_i\le-3\tau P/2\), so \(|Y|<1\).
An inward radial motion on an upper angle and an outward radial motion
on a lower angle both increase \(\Im W\):
\(\partial_{\log|y|}\Im W=-\cosh(\log|y|)\sin\theta\).
Thus all \(z_i\)'s stay on the interior sheet. Neither \(1-Z\)
nor a denominator \(1-z_iz_j\) vanishes. The canceled pair expression
has no remaining \(1-y_iy_j\) pole. This proves legality for the full
closed homotopy.

Normalize the angular Jacobian row \(i\) by \(y_i\). Its matrix is
\[
 M_{ij}=d_i\delta_{ij}+\frac{\lambda\tau}{2N}m_i p'_j,\qquad
 d_i=\ii+\lambda\left(-2\tau p'_i+\frac{\tau P}{2N}m'_i\right).
\]
Because \(m_ip'_i=0\), \(\det M=\prod_i d_i\); all factors are
nonzero and its absolute value is bounded by \(C^N\).
At a named \(q\), both the row and column of \(M\) are zero
off the diagonal and \(d_q=\ii\). Replacing angular column \(q\)
by the homotopy velocity multiplies the determinant by
\((-2\tau)/\ii=2\ii\tau\).
Cartan's formula on the closed torus gives
\[
 \partial_\lambda\int\psi U_{\lambda,s}
 =-2\ii\tau\sum_q\int(\partial_{\theta_q}\psi)U_{\lambda,s}.
\]
The unweighted integral is homotopy invariant. Integrating this identity
and separating the selected part proves \eqref{eq:FKS}.
\end{proof}

No equality is claimed between an individually weighted original chart
and the same weighted deformed chart. The current is indispensable.
All \(s\)-derivatives are taken on the entire fixed interval
\([0,1]\) before splitting the derivative density at \(\lambda_*\).
Thus no derivative of that split endpoint enters.

% PROV-INTERNAL: GENERIC_PAIR_TAIL.md; L2_PAIR_TAIL.md.
% RECONSTRUCTION: exact Schur identity, convolution and matching count.
\begin{proposition}[Ultra-high orders]\label{prop:ultrahigh}
For fixed \(\sstar\) and \(0\le j\le k\), there are constants
\(C_0,C_1,c>0\), independent of \(N,\eps\), such that
\[
 |T_N^{(j)}(\se)|
 \le C_j\eps^{-j-2}\left(\frac{C_1(H+1)}{\sqrt N}\right)^N.
\]
For a sufficiently large fixed \(C_H\),
the sum over even \(N\ge C_H(H+1)^2\) is \(o(\eps^A)\)
for every fixed real \(A\).
\end{proposition}
\begin{proof}
On the original \(c\eps\) disk, each root lies in the lower unit
half-disk. The identity
\[
 |1-zw|^2-|z-w|^2=(1-|z|^2)(1-|w|^2)+4\Im z\Im w
\]
and the Schur identity give
\(|\Pf A(z)|=\prod_{i<j}|(z_i-z_j)/(1-z_iz_j)|\le1\).
The product numerator is bounded by \(C^N\) and the two global
denominators below by \(c\eps\).
At the two simple inverse-dispersion branch angles,
\(|R(re^{\ii\theta};s)|^2\le C/(\eps+\dist(\theta,\{\pm\theta_B\}))\).
Its squared \(L^2\) norm is \(O(H+1)\). The \(y\)-pair kernel has a
majorant depending only on \(\theta_i+\theta_j\), with \(L^1\) norm
\(O(H+1)\). Circular convolution and Cauchy--Schwarz therefore bound
one matched weighted pair integral by \(C(H+1)^2\).

The \(y\)-Pfaffian has \(N!/[2^{N/2}(N/2)!]\) matchings.
Each matching factorizes after these majorants. Dividing by \(N!\)
and using \((N/2)!\ge(N/(2e))^{N/2}\) proves the undifferentiated
bound. Cauchy on the same disk adds \(j!(c\eps)^{-j}\).
Choose \(C_H\) so that the bracket is at most \(1/2\) at the threshold.
The summed bound is then
\(C_j\eps^{-j-2}2^{-C_H(H+1)^2}\), which has the stated decay.
\end{proof}

For \(N<C_H(H+1)^2\), denote the already differentiated,
angularly integrated current by
\begin{gather*}
 \mathcal J_{N,j}(\lambda)
   =-2\ii\tau\sum_q\int(\partial_{\theta_q}\psi)\,
              \partial_s^j U_{\lambda,s}\big|_{s=s_\eps},\\
 B_{N,j}=|K_N^{(j)}(s_\eps)|
       +\left|\int_0^{\lambda_*}\mathcal J_{N,j}(\lambda)\dd\lambda\right|.
\end{gather*}
The quantitative requirements are
\begin{align}
 \sum_N|F_N^{(j)}|
   &\le \exp(C_j\log^2(H+2)),\label{eq:reqF}\\
 \sum_N\left|\int_{\lambda_*}^1
               \mathcal J_{N,j}(\lambda)\dd\lambda\right|
   &\le C_j\eps^{-\alpha_0(j+2)},\label{eq:reqLS}\\
 B_{N,j}
   &\le C_j^N\eps^{-j-2}e^{-\kappa N^2}
             &&(N\ge D\sqrt H),\label{eq:reqhigh}\\
 B_{N,j}
   &\le \exp(C_jN\log(N+1))(H+N)^{C_j}
             &&(N_0+2\le N<D\sqrt H).\label{eq:reqinter}
\end{align}
The sums are over even tail orders \(N>N_0\) below \(C_H(H+1)^2\).
These are statements with fixed point and \(j\), followed by constants,
followed by a single sufficiently small \(\eps_0\), and only then all
orders in the indicated window.

% END SOURCE


% BEGIN SOURCE: sections/06_shifted_contour_F.tex
\section{The final selected contour}
\label{sec:F}
% PROV-INTERNAL: UPPER_SELECTED_MATCHING_FRESH_RECONSTRUCTION.md;
% GENERIC_UPPER_ANCHOR_AUDIT.md; UPPER_SELECTED_MATCHING_RECHECK.md.
\begin{proposition}\label{prop:F}
The bound \eqref{eq:reqF} holds for the final selected integral.
\end{proposition}
\begin{proof}
Telescope \(1-\psi\). A named selected index has \(p_q=1\), hence \(P\ge1\).
At \(\lambda=1\), this index is shifted inward by \(2\tau\).
The endpoint separation above gives
\(\Im W_q\ge\sinh(2\tau)\sin\delta\) at the radial center.
The interior root therefore stays in a strict subdisk, uniformly
on the selected support. It supplies a fixed \(Z\)-gap, and
\(\sum h_i\le-3\tau P/2\) supplies a fixed \(Y\)-gap.
Every true lower branch index is on \(m=1,p=0\) and has
\(\Im W\ge cP/N\), so its residue factor is at most \(C\sqrt N\).
Other roots are in a compact regular set or on the shifted upper
branch. The one-body product and pullback Jacobian cost
\(C^N N^{N/2}\).

These statements continue to the common disk \(|s-\se|\le c/N\):
the lower branch margin is \(c'/N\), compact roots move \(O(1/N)\),
and the fixed anchor absorbs the total compact-root inflation
\(\exp(CN(c/N))\) after the disk constant is made small.
The integral on this disk is holomorphic because the contour weights
are fixed in \(s\). Only the original undeformed representation needs
the smaller \(c\eps\) disk.

The assignments in the following absolute estimates are made
pointwise at each fixed \(s\). Their domains need not be holomorphic
in \(s\), and no assigned integral will be differentiated.
Let \(M\) be the number of indices outside the true lower branch core.
Split their compact \(z\)-images into a strict interior group and a
compact lower-arc group. On this disk, same-group compact Schur
pairs have modulus at most \(q<1\); cross compact pairs are at most
\(1+C/N\); branch/compact pairs are at most \(1+C/\sqrt N\);
branch/branch pairs are at most one.
Since the same-group count is at least \(M^2/4-M/2\),
\[
 |\Pf A(z)|\le
 \exp(CN+CM\sqrt N-aM^2)\le C_1^N e^{-a_1M^2}.
\]
The last inequality uses
\(CM\sqrt N\le(a/2)M^2+C^2N/(2a)\); the mixed inflation cannot be
dropped without this step.

Every \(y\)-pair with a true branch index is bounded. Indeed at
reciprocal angles its reflected upper partner has \(p=1\), and
the negative shift \(-2\tau\) dominates the possible positive
lower shift \(\tau P/(2N)\le\tau/2\).
For two compact angles write
\(d=\dist(\theta_i+\theta_j,2\pi\Z)\).
At exact reciprocity \(h_i+h_j\le0\) for every \(P/N\in[0,1]\);
smoothness gives \(h_i+h_j\le Ld\).
For \(d\le c\eps\) the negative base \(2\log r\) bounds the
modulus gap below by \(c'\eps\); for larger \(d\) the phase
bounds it below by \(c'd\). Thus
\[
 |A(y_i,y_j)|\le\frac C{\eps+d}.
\]
This majorant is independent of all other angles, including the
coupled scalar \(P\). Its matched two-angle integral is \(O(H+1)\).

Write the \(y\)-Pfaffian matrix as \(A_0+E\), where \(E\) has entries
only between compact indices and \(A_0\) has bounded entries.
Expand in an even compact subset of size \(2r\). The remaining
Pfaffian minor is at most \(C^N N^{(N-2r)/4}\) by Hadamard.
There are at most \(N^{3r}\) coarse choices of subset and matching;
their disjoint kernels have integral at most \((C(H+1))^r\).
As \(r\le M/2\), a sufficient bound is
\[
 \int|\Pf A(y)|\le C^N N^{N/4}(CN^3(H+1))^M.
 \]
Combine these bounds and \(N!\ge(N/e)^N\) to obtain, for each assignment,
\[
 |F_{N,\mathrm{assignment}}(s)|
 \le C^N N^{-N/4}(CN^3(H+1))^M e^{-a_1M^2}.
 \]
At most \(C^N\) assignments and a polynomial number of named indices
are absorbed by increasing \(C\). For \(N<C_H(H+1)^2\), completing
the square in \(M\) costs \(\exp(C\log^2(H+2))\).
This gives the same bound for the complete \(|F_N(s)|\) throughout
the common disk. Its fixed angular weights and the uniform integrable
majorant above make \(F_N\) holomorphic on a neighborhood of that
closed disk, after decreasing \(c\) once. Apply Cauchy to this complete
integral:
\[
 |F_N^{(j)}(\se)|\le j!\left(\frac Nc\right)^j
       \sup_{|s-\se|\le c/N}|F_N(s)|.
\]
No analytic dependence of the pointwise assignments is needed.
Since \(\sum_N C^N N^{-N/4+j}\) converges, this proves \eqref{eq:reqF}.
\end{proof}

% END SOURCE


% BEGIN SOURCE: sections/07_original_contour_K.tex
\section{The original contour and the all-branch region}
\label{sec:K}
% PROV-INTERNAL: NESTED_ANGULAR_PARTITION.md;
% SECOND_FAMILY_DESOURCE_SUPPORT_MATRIX.md, lines 21-31;
% SECOND_FAMILY_TAIL_GLOBAL_JET_AUDIT.md.
\begin{lemma}[Fixed support partition]\label{lem:partition}
The original \(K\)-integral has a finite smooth partition into an
all-lower-branch region and regions of the following types: a compact
left root with a modulus gap, a separated compact/true-branch pair,
and an all-compact right region. The partition uses at most \(C^N\)
labels, and every finite cutoff jet retains its named anchor and its
stated support.
Each named small-\(\lambda\) Stokes current has the same three
compact/branch types, with no all-branch sector.
\end{lemma}
\begin{proof}
Use \(u=\theta+\theta_B\) near the lower branch. Fix a small
\(\delta_B>0\) and a large fixed \(L\).
Take \(\chi_B=1\) for \(|u|\le\delta_B/2\), supported in
\(|u|\le\delta_B\). Take \(\chi_{\mathrm{in}}=1\) for
\(|u|\le r_0=\delta_B/(2L^2)\), supported in \(|u|\le2r_0\).
The deformation function \(m\) is one on a neighborhood of these
supports.
Split \(1=\prod\chi_B+(1-\prod\chi_B)\) and telescope the second
term. It has a named \(q\) with \(|u_q|\ge\delta_B/2\),
including on the support of every derivative of that named factor.
Expand the other coordinates with \(1=\chi_{\mathrm{in}}+(1-\chi_{\mathrm{in}})\).

To make the left/right labels periodic, choose an open arc \(J\)
whose closure lies inside the upper \(a=1\) plateau. Cut the circle
at a point of \(J\) and unwrap the angle on its complement. Near
\(-\theta_B\), use \(u=\theta+\theta_B\) and a smooth right step
which is zero for \(u\le-r_0/4\) and one for \(u\ge r_0/4\).
Away from this local transition keep its value constant on each
side, and return smoothly from one to zero inside \(J\).
This defines a smooth periodic right label; the left label is its
complement.

At the branch transition every derivative of the products with
\(1-\chi_{\mathrm{in}}\) vanishes, since \(\chi_{\mathrm{in}}=1\) there. The additional
transition inside \(J\) is also harmless: \(1-a\) vanishes there
with all its jets. Thus the full \(K\)-weight and every current
weight \(\partial_{\theta_q}\psi\), together with all their jets,
vanish whenever any coordinate lies in \(J\). For the named current
coordinate this follows from \(a'=0\); for every other coordinate
it follows from its factor \(1-a\). On the support of any relevant
weighted jet the labels therefore select the required side.
On the transition of \(\chi_{\mathrm{in}}\) they are constant. Every effective
compact label excludes a fixed branch neighborhood.
If a true-branch label occurs, the outer
\(q\) and a true-branch \(j\) satisfy
\(|b_j|\ge4|b_q|\) after choosing \(L\) large, unless a left root
already supplies a modulus gap.
With no branch label, all roots are compact; either one is left or
all are right. Finite product differentiation creates at most
\(N^{O(j)}\) terms and preserves the supports just described.

For the current fix its named index \(q_S\) in
\(\operatorname{supp}a'\). Split its two upper endpoint arcs;
this split can be smooth where the current weight vanishes.
An upper-left \(q_S\) gives a fixed root modulus gap.
For an upper-right \(q_S\), expand every other index with
\(\chi_{\mathrm{in}}+(1-\chi_{\mathrm{in}})\) and the compact
left/right labels. A compact left index again gives a modulus
gap. Otherwise, if the true-branch label set \(J_B\) is nonempty,
use \(q_S\) and the fixed index \(\min J_B\) as the mixed pair;
if it is empty all indices are compact right. These alternatives
are exhaustive, including when all \(N-1\) other indices are
true branch. The named upper index prevents an all-branch sector.

The slope of \(q_S\) is uniformly bounded on a preliminary
endpoint neighborhood separated from the upper branch, independently
of subsequently shrinking the slivers. Choose the nested branch
margin so that its limiting slope dominates this bound by more
than four. Lemma~\ref{lem:branch}'s deformed estimates preserve
\(|b_{q_S}/b_j|\le1/4\) after the final choice of \(\eps_0\).
The named index remains on \(p=1,m=0\), and every true branch
on \(p=0,m=1\), on neighborhoods of all cutoff-jet supports.
The vanishing on \(J\) proved above applies also to the current.
Thus no cutoff derivative creates a new support type.
\end{proof}

% PROV-INTERNAL: PAIR_CONSTANTS_COMPLEX_DISK_FRESH_AUDIT.md;
% GLOBAL_PAIR_CONSTANTS_INTERFACE.md; PAIR_COMPACTNESS_ENDPOINT_EQUALITY_AUDIT.md.
\begin{lemma}[Pair suppression at radial centers]\label{lem:contraction}
After the fixed branch margins are chosen, the endpoint slivers and
deformation strength can be chosen so that at \(s=s_\eps\), on
\(K^N\) and on each Stokes support with one named \(q\),
\[
 \left|\prod_{i<j}P_{ij}\right|\le C^Ne^{-\kappa N^2},
 \qquad \kappa>0 .
\]
This is uniform for \(0\le\lambda\le1\), all permitted occupancies
\(P/N\in[0,1]\), and sufficiently small \(\eps\).
More precisely, lower indices can be assigned to three fixed
covering groups with same-group bound \(q'<1\) and cross-group
bound \(1+\eta\), where \(\log(1+\eta)<-\log(q')/4\).
Pairs containing the named Stokes index are bounded by \(C_S\ge1\).
For every set \(E\) of at most \(j\) pairs, these bounds imply
\[
 \prod_{\substack{i<l\\(i,l)\notin E}}|P_{il}|
       \le C_j^N e^{-\kappa N^2}.
\]
Fixed pair jets are bounded in regular branch coordinates and
compact angular coordinates; no uniform bound on raw branch-angle
jets is asserted.
\end{lemma}
\begin{proof}
\noindent\textbf{Fixed groups and limiting equality loci (rc4 repair).}
Write \(b=\theta_B\), \(y_B=e^{-\ii b}\), and
\(t_0=S_*+1-\sqrt{(S_*+1)^2-1}\in(0,1)\).
Choose a fixed \(w<\min(b,\pi-b)/3\), small enough for the
regular lower-branch chart, with \([-b-w,-b+w]\) inside the
\(m=1\) neighborhood. Choose \(0<\ell_L,\ell_R<w\).
The closed limiting groups are
\[
 B=[-b-w,-b+w],\quad L=[-\pi,-b-\ell_L],\quad
 R=[-b+\ell_R,0].
\]
They cover the lower semicircle with overlaps. Assign each overlap
by any fixed rule. Add the upper-left sliver \([\pi-2\delta,\pi]\)
to \(L\) and the upper-right sliver \([0,2\delta]\) to \(R\),
identifying \(\pi\) and \(-\pi\). The named current index lies
in one of the two transition arcs. All cutoff jets retain these
closed supports. The branch group contains neither endpoint;
otherwise \(|P(y_B,1)|=|P(y_B,-1)|=1\) would prevent strictness.

At the limit \(W_0(\theta)=S_*-\cos\theta\) and
\(W_0+1\ge S_*>0\). On \([-\pi,-b]\) the roots are real
in \([t_0,1]\); on \([-b,0]\) they are \(e^{-\ii v}\),
\(0\le v\le b\). The identity
\[
 |1-zw|^2-|z-w|^2
 =(1-|z|^2)(1-|w|^2)+4\Im z\Im w
\]
and its unit-circle sine version give \(|P|\le1\).
Apart from the removable configurations described below, equality
requires an endpoint in the \(y\)-pair and equality in the
\(z\)-factor. Its complete loci, with either ordering, are
\(y_i=1,\ \theta_j\in[-\pi,-b]\), and
\(y_i=-1,\ \theta_j\in[-b,0]\). Common-endpoint pairs have
canceled value zero. No equality locus lies in a same-group square.
On \(L^2\) the roots stay in a strict real subinterval of \((0,1)\);
on \(R^2\) they stay on a compact lower arc away from the real axis.
On \(B^2\) the \(y\)-factor is bounded by
\[
 \frac{\sin w}{\min\{\sin(b-w),\sin(b+w)\}}
 \le\frac1{2\cos w}<1.
\]
The canceled extension below covers its double branch. Thus
the three compact same-group suprema have a maximum \(q<1\).
Cross-group suprema may equal one.

\par\smallskip\noindent\textbf{Local continuation and the complete pair.}
Use thin simply connected neighborhoods of the compact real
\(W_0>1\) intervals and of the compact intervals inside \((-1,1)\),
avoiding both branch points. On the former choose the root real
in \((0,1)\); on the latter choose the continuation of
\(e^{-\ii\arccos W_0}\). These local roots may leave the unit disk.
Near the lower branch use \(\phi=\arccos W\) on the
fourth-quadrant sheet for \(\Im W\ge0\).
The implicit-function theorem gives the even analytic inverse
\(y_s(\phi)=G_s(\phi^2)\) near \((s_*,0)\), since
\(\partial_yW(y_B;s_*)\ne0\). In this neighborhood
\[
 y_s(\phi_i)-y_s(\phi_l)
  =(\phi_i-\phi_l)(\phi_i+\phi_l)H_s(\phi_i,\phi_l),\qquad
 1-e^{-\ii(\phi_i+\phi_l)}=(\phi_i+\phi_l)J(\phi_i+\phi_l),
\]
where \(J(0)=\ii\). Substitution into \eqref{eq:pair} gives
\(P_{il}=(\phi_i-\phi_l)^2A_s\), with analytic bounded fixed
jets and \(A_{s_*}(0,0)=-1/(4\sin^2b)\).
This resolves the whole local reciprocal-root hypersurface,
without dividing by a collision zero.

At the limiting lower support \(1-z_i z_l=0\) forces both
roots to be \(1\) or both \(-1\). The latter is excluded by
\(W_0+1\ge S_*\); the former is the resolved lower double branch.
On the closed compact complement of a smaller double-branch
neighborhood, the denominator therefore has a positive minimum
\(\mu\). Shrink its neighborhoods to retain \(\mu/2\).
The regular and rational expressions agree on overlaps by
\eqref{eq:pair}. This finite cover supplies continuous complete
pairs and bounded fixed jets in regular branch and compact
coordinates. No globally glued square-root domain is needed.
The apparent \(1-y_i y_l\) pole was already canceled; the upper
inverse branch, which would make a genuine pole with the lower
branch, is excluded by the selector plateau.

For a named right limiting endpoint root \(\xi=e^{-\ii b}\),
\(|1-\xi z|=|e^{\ii b}-z|\ge\sin b\) for every limiting lower
root. For the left endpoint it is at least \(1-t_0\).
Thus a positive named margin
\(\mu_S=\min(\sin b,1-t_0)\) persists, with \(\mu_S/2\)
on sufficiently small neighborhoods. Boundedness of the canceled
numerator supplies \(C_S\ge1\); no named contraction is needed.

\par\smallskip\noindent\textbf{Uniform support transfer and reserved slack.}
Project each sliver to its endpoint and keep each lower angle
fixed; denote the projected angle by \(\bar\theta_i\).
For \(\rho=P/N\in[0,1]\) and \(0\le\lambda\le1\),
\[
 v_i=\log|y_i|=-c_0\eps+\lambda\tau(-2p_i+\rho m_i/2),
 \qquad |v_i|\le c_0\eps+5\tau/2.
\]
Let \(L_S\) bound \(|S'|\) on a fixed parameter neighborhood.
Uniformly in this entire parameter box, hence for every actual
coupled occupancy and every \(N\),
\[
 |y_i-e^{\ii\bar\theta_i}|+|W_i-W_0(\bar\theta_i)|
 \le C(\eps+\tau+\delta)+L_S|s-s_\eps|.
\]
Compact-root changes are Lipschitz. Near the branch, the factorization
\(\phi=\sqrt{2(1-W)}h(1-W)\), with analytic \(h\), gives a
one-half Holder bound on the chosen closed half-plane.
The regular pair is Lipschitz in \((s,\phi_i,\phi_l)\);
on the compact complement the denominator margin gives the same
property in \((y_i,y_l,z_i,z_l)\). A finite cover therefore gives
one uniform complete-pair modulus, bounded by
\(C\sqrt{\eps+\tau+\delta+|s-s_\eps|}\) in these neighborhoods.
It applies whenever the branch-near roots stay on the designated
sheet; the two disk proofs below verify this condition separately.

Choose \(q<q_1<q'<1\), then \(\eta>0\) with
\(\log(1+\eta)<-\log(q')/4\).
Choose the slivers and \(\tau\), and a preliminary \(\eps_0\),
so the radial-center modulus is smaller than both \(q_1-q\)
and \(\eta/2\), while retaining the denominator neighborhoods.
This gives same-group bound \(q_1\) and cross-group bound
\(1+\eta/2\); reserve \(q'-q_1\) and \(\eta/2\) for the disks.
The real radial motions increase \(\Im W\), so every radial-center
root is interior. This transfer includes the slivers in the
same-group estimate as well as the cross estimate.

For clarity, the limiting branch/upper-sliver moduli for small
\(a>0\) are
\[
 \frac{\sin((b+a)/2)}{\sin((b-a)/2)}
   =1+a\cot(b/2)+O(a^2),\qquad
 \frac{\cos((b-a)/2)}{\cos((b+a)/2)}
   =1+a\tan(b/2)+O(a^2).
\]
Holding a sliver fixed while moving a compact-group boundary
toward the branch can consequently make an \(RR\) or \(LL\)
pair exceed one. Such a choice does not obey the required
margin--slack--sliver order; it is a real geometric excess.

Let \(n\) be the number of lower variables: \(n=N\) for \(K\)
and \(n=N-1\) for a Stokes term. Assign each to one covering group.
After at most \(j\) pairs have been removed, the remaining
same-group count is at least \(n^2/6-n/2-j\), while the
remaining cross-group count is at most \(n^2/3\). Apply the
individual upper bounds directly to the remaining product:
\[
 (q')^{n^2/6-n/2-j}(1+\eta)^{n^2/3}
       \le(q')^{n^2/12-n/2-j}.
\]
This remains valid if the lower bound on the count is negative.
It involves no division by a removed factor, which may vanish.
The named upper compact
\(q\) has at most \(N-1\) additional pairs, each uniformly
bounded: its \(z_q\) is separated from the inverse of a lower
root, including the branch root. These pairs cost \(C_S^{N-1}\).
For \(n=N-1\) the displayed exponent is
\(N^2/12-2N/3+7/12-j\); the linear correction is included in
\(C_j^N\). Thus \(\kappa=-\log(q')/12\) works for both objects.
The canceled double-branch expression is analytic in the regular
coordinates, and the other pair denominators are separated on the
fixed compact sets. This also proves the stated fixed-jet bounds.
\end{proof}

\begin{lemma}[The original parameter disk]\label{lem:originaldisk}
For a sufficiently small fixed \(c>0\), the original contour
density and all \(K\) and Stokes transition densities are holomorphic
on \(|s-s_\eps|\le c\eps\). On the \(K\) and named Stokes
supports, the following estimates hold:
\[
 |1-Y|,\ |1-Z|\ge c'\eps,\qquad
 \left|\prod_{i<j}P_{ij}\right|\le C^Ne^{-\kappa' N^2}.
\]
The integrated one-body majorants and contour Jacobian cost
at most \(C^N\), uniformly in \(N\) and \(0\le\lambda\le1\).
\end{lemma}
\begin{proof}
The undeformed dispersion has \(\Im W\ge c_2\eps\).
Every permitted radial motion increases this quantity, as in
Proposition~\ref{prop:selector}; reducing \(c\) preserves a
fixed fraction of the margin on the disk. Each root is then
in the lower unit half-disk with attenuation at least \(c_3\eps\).
Thus the product gaps are bounded below by a constant times
\(\min(1,N\eps)\), which is at least a constant times \(\eps\).
The contour \(Y\) is independent of \(s\).

Near the lower branch put \(t=\lambda P/N\),
\(E=\eps+t\), and \(\delta=|s-\se|\le c\eps\).
On its true plateau, write
\(y=e^{v+\ii\vartheta}\), where
\(v=-c_0\eps+\tau t/2\) and \(\vartheta=-\theta_B+u\).
The exact identities
\[
 \Re W=\Re S-\cosh v\cos\vartheta,\qquad
 \Im W=\Im S-\sinh v\sin\vartheta
\]
give, with \(D_B=1-W\) and \(a_B=\sin\theta_B>0\),
\[
 |\Re D_B-a_Bu|
 \le C(u^2+\eps^2+t^2+\delta),\qquad
 \Im W\ge c_bE-C_S\delta\ge(c_b/2)E.
\]
Here the last inequality holds after reducing the disk constant
\(c\). All constants depend only on the fixed point and margins.
In particular, the real branch-center displacement on this disk is
\(O(\eps^2+t^2+\delta)\): the parameter perturbation cannot be
omitted from the radial-center estimate.

Choose a fixed branch radius \(h\) and a fixed \(K\ge1\).
If \(|u|\le KE\), the imaginary margin gives
\(|D_B|\ge c_b(|u|+\eps)/(2(K+1))\).
If \(KE<|u|\le h\), then
\(u^2\le h|u|\), \(\eps^2+t^2\le E^2\le h|u|/K^2\),
and \(\delta\le cE\le c|u|/K\).
Choose \(h,K,c\) so that
\(C(h+h/K^2+c/K)\le a_B/2\), in addition to the previous margins.
Then \(|\Re D_B|\ge a_B|u|/2\) in this second region.
These choices are independent of \(N,\eps,\lambda,P\).
Since \(W+1\) stays separated from zero on the fixed branch arc
and \(R=z/\sqrt{W^2-1}\), both regions yield
\[
 |R(y;s)|\le \frac{C}{\sqrt{|u|+\eps}}
\]
uniformly on the full disk. The upper branch is absent from
these supports; the remaining compact factors are bounded.
The displayed majorant does not contain the coupled quantity
\(P\), and its integral on \(|u|\le h\) is at most \(4C\sqrt h\).
Thus its product can be integrated with cost \(C^N\) without
treating \(P\) as constant. The inverse-product numerator
and the Jacobian in Proposition~\ref{prop:selector} also cost \(C^N\).

For an explicit disk choice, let \(b_0\eps\) be the
undeformed imaginary margin and let \(L_S\) be as in
Lemma~\ref{lem:contraction}. Choose \(L_Sc\le b_0/4\).
The monotone radial motions then leave \(\Im W\ge3b_0\eps/4\)
throughout the disk. Along each disk radius,
\(\partial_W\phi=-1/\sin\phi\) gives branch-coordinate motion
\(O(c\sqrt\eps)\); compact roots move by \(O(c\eps)\).
The finite continuation cover and the regular canceled expression
of Lemma~\ref{lem:contraction} give a uniform complete-pair change
\(O(\sqrt\eps)\). Reduce the final \(\eps_0\) to make it
smaller than both reserved margins \(q'-q_1\) and \(\eta/2\).
Thus same pairs are at most \(q'\), cross pairs at most
\(1+\eta\), and named pairs are uniformly bounded.
The limiting cross supremum need not be strict. The direct
remaining-pair count proves the estimate with the same
\(\kappa=-\log(q')/12\), or a smaller common exponent.
These pair estimates concern the stated \(K\) and current
supports, which exclude the upper inverse branch. They do not
assert a uniform pair bound on the unrestricted original torus;
that contour is handled separately in Proposition~\ref{prop:ultrahigh}.
No large-\(\lambda\) estimate is used here.
\end{proof}

% PROV-INTERNAL: GENERIC_BRANCH_REAL_PHASE_GEOMETRY.md;
% ALL_B_LIE_NUMERATOR_DESOURCE.md.
\begin{lemma}[Branch geometry]\label{lem:branch}
On a sufficiently small original lower-branch arc,
\[
 1-W=a_Bu-\ii b_B\eps+O(u^2+\eps|u|+\eps^2),
 \quad a_B=\sin\theta_B>0,\quad b_B=2\sin\theta_*-c_0a_B>0.
\]
With \(\phi=\arccos W\) on the fourth-quadrant sheet,
\begin{align*}
 |\phi'(u)|&\asymp(|u|+\eps)^{-1/2},&
 \int_B|\phi'(u)|\dd u&\le C,\\
 -\Im\phi(u)&\ge c\sqrt{|u|+\eps}\quad(u\le0),&
 x'(u)&\asymp(u+\eps)^{-1/2}\quad(u\ge0),
\end{align*}
where \(x=\Re\phi\).
For \(u\ge C_0\eps\), \(-x''(u)\ge c u^{-3/2}\).
If \(M/\eps\to\infty\) uniformly and \(0\le m\le M/2\),
\[
 x'(m)-x'(M)\ge c(m+\eps)^{-1/2},\qquad
 \frac{|\phi'(m)\phi'(M)|}{x'(m)-x'(M)}\le C M^{-1/2}.
\]
If \(M/2<m<M\) and \(m\gg\eps\), the latter ratio is at most
\(C\sqrt M/(M-m)\).

On a true-branch plateau of the current put
\(t=\lambda P/N\), \(E=\eps+t\). Uniformly for
\(0\le\lambda\le\lambda_*\) and \(|u|\le2r_0\),
\begin{align*}
 |\phi_u|&\asymp(|u|+E)^{-1/2},\\
 -\Im\phi&\ge c\sqrt{|u|+E}\quad(u\le0),&
 \Re\phi_u&\ge\tfrac23|\phi_u|\quad(u\ge0).
\end{align*}
The constants may depend on the fixed \(\tau\), but not on
\(N,\eps,\lambda,P/N\). Here the margins and then \(\eps_0\)
are small enough, and \(P\) is fixed during a plateau variation.
\end{lemma}
\begin{proof}
The displayed normal form follows by Taylor expansion of \(S\) and
the fixed-radius \(y\). Put \(D=1-W\).
The exact factorization \(\phi=\sqrt{2D}\,h(D)\), \(h(0)=1\),
with analytic \(h\), yields the first estimates and an integrable
branch length. On \(u\ge C_0\eps\), the leading terms
\(x'(u)=\sqrt{a_B/2}\,u^{-1/2}(1+O(\eps/u+u))\) and
\(-x''(u)\asymp u^{-3/2}\) have the stated signs after increasing
\(C_0\) and shortening the arc.
If \(m\ge C_0\eps\) and \(m\le M/2\), the relative error in
\(x'\) can be made smaller than a fixed fraction of
\(1-1/\sqrt2\); this gives the strong difference estimate.
If \(m<C_0\eps\), its slope is at least \(c\eps^{-1/2}\),
whereas the slope at \(M\) is \(o(\eps^{-1/2})\).
For comparable \(m,M\), integrate the bound on \(-x''\).
The sector containing \(\phi'\) also has width less than \(\pi\);
integrating in \(u\) gives
\(|\phi(u)-\phi(v)|\ge c|u-v|\).

For the current, write \(y=e^{v+\ii\vartheta}\),
\(v=-c_0\eps+\tau t/2\), \(\vartheta=-\theta_B+u\).
Put \(Q=\eps+\tau t\), comparable to \(E\) for fixed \(\tau>0\).
The exact real and imaginary identities in
Lemma~\ref{lem:originaldisk}, at the radial center, give
\begin{gather*}
 D=1-W=a_Bu-\ii\beta+R,\qquad cQ\le\beta\le CQ,\\
 R\in\R,\qquad |R|\le C(u^2+Q^2),\qquad
 D_u=a_B+O(|u|+Q).
\end{gather*}
Positivity of \(\beta\) uses \(b_B>0\) and the positive
\(\tau t\) coefficient. The constants in these \(Q\)-bounds
can be chosen before shrinking \(\tau\). The model
\(a_Bu-\ii\beta\) has modulus comparable to \(|u|+Q\), and
its relative perturbation is \(O(r_0+Q)\). Its square root for \(u\le0\) has argument
in \([-\pi/2,-\pi/4]\); its inverse square root for \(u\ge0\)
has argument in \([0,\pi/4]\). Using the exact factorization
\(\phi=\sqrt{2D}\,h(D)\) and its derivative proves the asserted
cone and slope bounds after shrinking the fixed margins and
\(\eps_0\). This argument allows the \(O(E^2)\) real displacement
of the branch center; it does not require \(\Re D\le0\) at \(u=0\).
In particular, since all \(|z_i|\le1\), for a negative branch
coordinate
\[
 \frac{|\phi'_v|}{|1-Z|}
 \le\frac{|\phi'_v|}{1-|Z|}
 \le\frac C{|u_v|+\eps+\lambda P/N}.
\]
\end{proof}

% PROV-INTERNAL: NONRESONANT_BRANCH_EXPONENTIAL_CUTOFF.md;
% NONRESONANT_BRANCH_HOSTILE_AUDIT.md; TWO_NONZERO_G3_FRESH_COAREA_PROOF.md.
\begin{proposition}[Nonresonant microcore]\label{prop:micro}
For \(N_0+2\le N<D\sqrt H\), put
\[
 b_N=c_*e^{-(A+4)N}/N,\qquad
 \psi_N=\prod_i\chi(u_i/b_N),
\]
where \(\chi=1\) on \([-1/2,1/2]\), supported on \([-1,1]\), and
\(c_*>0\) is fixed and small. The all-\(B\) microcore contribution
obeys, for every \(j\le k\),
\[
 |\partial_s^j I_{\mathrm{micro},N}|
 \le \exp(C_jN^2)\,b_N^{(N^2-1)/2}.
\]
\end{proposition}
\begin{proof}
Uniformly in this window,
\(N\eps e^{AN}\to0\) and \(\eps/b_N^m\to0\) for every fixed \(m\).
Thus on the microcore and all its cutoff jets,
\[
 |Y-\xi^N|\le CN(\eps+b_N)\le\tfrac12e^{-AN},
 \qquad |1-Y|\ge\tfrac12e^{-AN}.
\]
The one-phase field \(V_i=a_i=S'/W_{u_i}\) freezes each regular
\(\phi_i\), hence \(Z\). In these coordinates the numerator is
analytic times \(\Delta(\phi)^2\).
Repeated Lie differentiation leaves the latter polynomial unchanged
and leaves the \(Z\)-kernel simple. The \(Y\)-power is at most
\(j+1\); its gap and the \(b_N^{-j}\) cutoff cost are
\(\exp(O_j(N))\). Analytic products have bound \(\exp(O_j(N^2))\).

The branch image lies in a fixed cone of width less than \(\pi\)
and has radius \(C\sqrt{b_N}\). Decreasing \(c_*\) ensures
\(N\sqrt{b_N}\) is small, so
\(|1-Z|\ge c\sum_i|\phi_i|\).
The square-root normal form gives, directly on both sides of zero,
\(|\phi|\ge c\sqrt{|u|+\eps}\) and
\(|\phi'|\le C(|u|+\eps)^{-1/2}\). Substituting
\(w=\sqrt{|u|+\eps}\) therefore gives
\[
 \int_{\mathrm{core}}e^{-t|\phi|}|\dd\phi|
       \le4C\int_{\sqrt\eps}^{\sqrt{b_N+\eps}}e^{-ctw}\dd w
       \le C\min(\sqrt{b_N},t^{-1}).
\]
Using \(1/\sum|\phi_i|=\int_0^\infty e^{-t\sum|\phi_i|}\dd t\)
therefore costs \(C^N b_N^{(N-1)/2}\).
The Vandermonde costs at most
\(C^{N(N-1)}b_N^{N(N-1)/2}\). These estimates prove the claim.
The logarithm of the displayed majorant is
\[
 -\frac{A+4}{2}N^3-\frac12N^2\log N+O_j(N^2).
\]
In particular it is at most \(-(A+4)N^3/2+O_j(N^2)\).
Thus the sum of these bounds over all \(N\ge N_0+2\) is finite.
\end{proof}

% PROV-INTERNAL: TWO_NONZERO_G3_FRESH_COAREA_PROOF.md;
% ALL_B_LIE_NUMERATOR_DESOURCE.md; NONRESONANT_BRANCH_HOSTILE_AUDIT.md.
\begin{proposition}[The all-branch complement]\label{prop:allB}
The remaining original all-\(B\) contribution, for
\(N_0+2\le N<D\sqrt H\), is bounded through every \(j\le k\) by
\[
 \exp(C_jN\log(N+1))(H+N)^{C_j}.
\]
\end{proposition}
\begin{proof}
Telescope \(1-\psi_N\); each term has a named anchor \(a\) with
\(|u_a|\ge b_N/2\), including all cutoff jets.
Let \(t=N^{-1}\sum u_i\), \(\rho^2=\sum(u_i-t)^2\), and split
smoothly at \(\rho_N=e^{-BN}\), using a function of
\(\rho^2/\rho_N^2\). Put
\[
 \Gamma_* = \sup_{\substack{n\in\mathbb N\\n\ge N_0+2}}
                   \frac{\log(16n/c_*)}{n}<\infty.
\]
Choose \(B>A+4+\max\{1,\Gamma_*\}\), increasing it as needed
after the finite-jet constants have been fixed. Then
\[
 2\rho_N\le b_N/8\qquad\hbox{for every }N\ge N_0+2.
\]
This is a support inclusion for all orders, not an asymptotic
comparison repaired by enlarging finitely many estimate constants.

On the far support \(\rho\ge\rho_N\), the maximal separation
is at least \(\rho_N/\sqrt N\). The rational two-phase Lie
calculus costs \(\exp(O_j(N\log N))\), including every exponential
cutoff jet. The undifferentiated pairs retain
\(\exp(-\kappa N^2+O_j(N))\).
It remains to integrate two simple global kernels against
\(\prod_i|\phi'_i|\dd u_i\).

If all coordinates are nonnegative, their maximum \(M\ge b_N/2\).
Select the actual minimum \(m\) and maximum \(M\) only now,
after differentiation. For fixed other variables the real phase
map is
\[
 (m,M)\longmapsto(F,G)
    =(m+M+\mathrm{const},\,x(m)+x(M)+\mathrm{const}).
\]
Its Jacobian \(x'(m)-x'(M)>0\) makes it injective on each fixed-\(F\)
extreme-index cell. Lemma~\ref{lem:branch} bounds the ratio of its
two branch measures to that Jacobian by \(\exp(O(N))\):
use \(M\ge b_N/2\) for \(m\le M/2\), and
\(M-m\ge\rho_N/\sqrt N\) otherwise.
Only \(m,M\) vary here; all other angles are fixed. Each unwrapped
phase-image interval has bounded length, independent of \(N\), since
\(m,M\) and \(x(m),x(M)\) range over fixed bounded intervals.
The additive constants translate these intervals without changing
their lengths. Thus each simple periodic phase integral costs
\(O(H+N)\). The other branch lengths cost \(C^N\).

For negative coordinates choose \(\delta_N=e^{-C_1N}\) with
\(C_1>A+4+\max\{1,\Gamma_*\}\). Then \(\delta_N\le b_N/16\)
for every \(N\ge N_0+2\). If some negative \(|u_v|\ge\delta_N\), then
\(|1-Z|\ge e^{-O(N)}\); integrate the remaining \(Y\)-kernel
in the named anchor, whose measure is at most \(e^{O(N)}\dd u_a\).
If all negative magnitudes are smaller than \(\delta_N\), the
named anchor is positive. Choose one negative \(v\); the bound
\[
 \frac{|\phi'(u_v)|}{|1-Z|}\le\frac C{|u_v|+\eps}
\]
has integral \(O(H+N)\). The positive anchor integrates the
\(Y\)-kernel with cost \(e^{O(N)}(H+N)\).
The scale inclusion already covers every finite order; any
remaining finite-order constants only enlarge the numerical bound.
This proves the far bound
\(\exp(-\kappa N^2+O_j(N\log N))(H+N)^2\).

On the near support, \(|u_i-t|\le\rho\le2\rho_N\le b_N/8\).
The anchor \(|u_a|\ge b_N/2\) gives \(|t|\ge3b_N/8\), so
all coordinates have the same sign and \(|u_i|\ge b_N/4\).
The pointwise branch-measure product has the bound
\[
 \prod_i|\phi'_i|\le(Cb_N^{-1/2})^N
 =\exp\left(\frac{A+4}{2}N^2+
             \frac N2\log(N/c_*)+O(N)\right).
\]
Segments joining two coordinates stay on this same-sign support,
so each collision difference is at most \(Cb_N^{-1/2}\rho\).
The collision coefficient and the branch measures together are
therefore bounded by \(C^{N^2}b_N^{-N^2/2}\), whose logarithm is
\[
 \frac{A+4}{2}N^3+\frac{N^2}{2}\log(N/c_*)+O(N^2).
\]
Fixed-\(j\) coefficient and derivative costs can be included in
the lower-order \(O_j(N^2)\) term. Its constant is fixed before
\(B\). These bounds are determined
on the \(B\)-independent anchor neighborhood; the positive cubic
coefficient is independent of \(B\). Lie and cutoff jets lose at most
\(2j+1\) powers of \(\rho\).
The shape measure contributes \(\rho^{N-2}\).
For negative \(t\), the \(Z\)-gap is \(e^{-O(N)}\) and one mean
phase integral suffices.
For positive \(t\), strict concavity and zero first shape derivative
give \(-\partial_\rho\Re\sum\phi_i\ge c\rho\), while the \(Y\)-phase
is \(Nt+\mathrm{const}\). Their real Jacobian is at least \(cN\rho\).
After coarea the conservative remaining power is
\[
 L=N^2-2j-4\ge N^2-2k-4\ge4N_0+2>0.
\]
Bound this power by \((2\rho_N)^L\) before integrating the phase
kernels. At fixed shape direction, \(F=Nt+\mathrm{const}\) fixes
\(t\), and strict monotonicity of \(G\) in \(\rho>0\) gives
multiplicity one for the unwrapped map. Each phase lies in an
interval of length at most \(CN\): the mean angle ranges over a
fixed interval and each individual branch phase is bounded.
The periodic-kernel estimate in Appendix~\ref{app:lie} therefore
costs at most \(CN^2(H+N)^2\) for the two phases, rather than a
single-period bound. For negative \(t\), the sole mean phase
instead costs at most \(CN(H+N)\), with the separated \(Z\)-kernel
already in the preceding estimate. These polynomial factors are
absorbed in \(e^{C_jN\log(N+1)}\). The coarea variable is not
integrated a second time. The bound is at most
\[
 \exp\left(\frac{A+4}{2}N^3+\frac{N^2}{2}\log N+C_jN^2\right)
 (2\rho_N)^L(H+N)^2.
\]
Its logarithm, apart from the simple-kernel factor, is at most
\(((A+4)/2-B)N^3+\tfrac12N^2\log N+O_j(N^2)+B(2j+4)N\).
The already required \(B>A+4\) makes this tend to minus infinity.
Finitely many remaining \(N\)'s cost a fixed, possibly very large,
\(\eps\)-independent constant.
At fixed \(\eps>0\) the same positive collision power makes the
full-equality puncture flux vanish at every preceding Lie stage.
Thus the identities used on the punctured charts have no residual
boundary term. The near and far bounds prove the proposition.
\end{proof}

% END SOURCE


% BEGIN SOURCE: sections/08_stokes_S.tex
\section{Compact and mixed sectors; Stokes estimates}
\label{sec:sectors}
% PROV-INTERNAL: GENERIC_MIXED_STOKES_REDERIVATION.md;
% SMALL_LAMBDA_LIE_ATLAS_AUDIT.md; NESTED_ANGULAR_PARTITION.md.
\begin{proposition}[Separated mixed and left sectors]\label{prop:mixed}
For \(N_0+2\le N<D\sqrt H\), the non-all-\(B\) original \(K\)
sectors with a compact left root or a separated compact/branch
pair, and the corresponding small-\(\lambda\) Stokes sectors,
obey \eqref{eq:reqinter}. The differentiated current density
obeys the per-order bound uniformly for \(0\le\lambda\le\lambda_*\).
\end{proposition}
\begin{proof}
Use hybrid coordinates consisting of regular \(\phi_i\) for true
branch indices \(J\) and original angles for the compact indices.
All coefficient jets are bounded in these coordinates.
For Stokes the named \(q\) lies on \(p=1,m=0\), and every true
branch lies on \(m=1,p=0\); their angle variations preserve both
\(P=\sum p_i\) and \(\sum m_i\) exactly.
For original \(K\) the deformation is zero, so these occupancy
constraints are unnecessary.

Put \(a_i=(\partial_s\phi_i)/b_i\),
\(b_i=\partial_{\theta_i}\phi_i\),
\(A=\sum_{i\in J}a_i\), and
\(D=\sum_{i\notin J}\partial_s\phi_i\).
With \(j=\min J\) fixed for the labelled term, let
\[
 X=\frac{D+b_qA}{b_j-b_q},\qquad
 V_i=a_i\ (i\in J\setminus\{j\}),\quad
 V_j=a_j+X,\quad V_q=-A-X .
 \]
All other velocities vanish. Direct substitution freezes \(Y,Z\)
and every unselected branch phase. The residual branch velocity is
\[
 (\partial_s-V\cdot\nabla)\phi_j
       =-\frac{D+b_qA}{1-b_q/b_j}.
\]
The support choice gives \(|b_q/b_j|\le1/4\).
Since \(1/b_j\) is analytic and vanishes in its regular branch
coordinate, this residual, its finite jets, and the compact
velocity have only polynomial \(N\)-cost, uniformly in \(\eps,\lambda\).
Here the hybrid top form can be written explicitly. Up to the
unchanged factor \((Z^{-1}+Y^{-1})/[N!(1-Z)(1-Y)]\prod_{i<l}P_{il}\),
its measure is
\[
 \prod_{i\in J}\left[-\frac{z_i}{\pi(1-y_i^{-2})}\dd\phi_i\right]
 \prod_{i\notin J}\left[\frac{R_i y_i d_i}{2\pi\ii}\dd\theta_i\right].
\]
Indeed \(d_i=\ii\) on every true branch plateau, the angular
determinant is \(\prod d_i\), and the hybrid change of variables
has triangular determinant \(\prod_{i\in J}b_i\).
This retains the full Jacobian of the coupled contour.

For Stokes, \(V\cdot\nabla P=V\cdot\nabla\sum m_i=0\) identically
on an open neighborhood of the support. The same is true for
all parameter derivatives of this identity. Thus the coupled
occupancy introduces no hidden branch motion at higher order.
The actual current multiplier is
\(-2\ii\tau\,\partial_q\psi\); its finite derivatives, and the
derivatives of the nested partition, stay on these plateaus.
Use the pullback recurrence \eqref{eq:pullbackLie} in the hybrid
coordinates. Its residual coefficients have no difference pole,
because \(1-b_q/b_j\) is separated from zero. Each fixed operation
changes finitely many analytic pair or one-body factors and
introduces at most polynomial \(N\)-cost. The remaining pair
factors therefore retain
\(C_j^N N^{C_j}e^{-\kappa N^2}\).

For positive branch coordinates, Lemma~\ref{lem:branch}, including
its current version, gives \(\Re b_j\ge(2/3)|b_j|\).
The compact slope is dominated by the true branch slope; for a
Stokes upper-right \(q\) it has the opposite sign in the limit.
The phase map
\((\theta_q,u_j)\mapsto(\sum\theta_i,\Re\sum\phi_i)\)
has Jacobian \(\Re(b_j-b_q)\ge(2/3-1/4)|b_j|\)
and is injective on fixed first-phase slices.
It absorbs the selected branch measure, and each remaining phase
kernel integrates with cost \(O(H+\log N)\).
Every unselected branch measure has bounded integral.

If a branch coordinate \(u_v<0\), its attenuation and branch measure
give instead
\[
 \frac{|\phi'_v|}{|1-Z|}
 \le\frac C{|u_v|+\eps+\lambda P/N}.
\]
Integrate this factor in \(u_v\), and the \(Y\)-phase in the compact
\(q\). Both cost at most \(C(H+\log N)\).
The phase statements use variables on true plateaus; therefore \(P\)
does not vary during the selected integrations.
These sign divisions are performed only on the final absolute integral.

If a compact left root occurs, \(Z\) has a fixed modulus gap.
Set \(V_i=a_i\) for branches and \(V_q=-A\), preserving \(Y\)
and freezing every branch root. Differentiating the separated
\(Z\)-kernel costs only polynomial \(N\), and \(q\) integrates
the simple \(Y\)-kernel. A left \(q\) itself is allowed.
The fixed support partition exhausts the asserted cases.
\end{proof}

% PROV-INTERNAL: COMPACT_R_RATIONAL_ATLAS_PROOF.md;
% COMPACT_R_FRESH_KTH_COAREA_PROOF.md;
% COMPACT_R_STOKES_EQUALITY_DESOURCE.md.
\begin{proposition}[Compact right sectors]\label{prop:compactR}
The all-compact right original \(K\) and small-\(\lambda\) Stokes
sectors satisfy the same per-order bound in \eqref{eq:reqinter}.
\end{proposition}
\begin{proof}
On a fixed compact right interval \(R\), separated from both branch
points, put \(g(\theta)=\arccos(S_*-\cos\theta)\). Direct differentiation
gives
\[
 g''(\theta)=-\frac{S_*(1-\cos\theta\cos g(\theta))}{\sin^3g(\theta)}
            \le-c_R<0 .
\]
The coupled contour is smooth in the angles and analytic in \(s\).
Its fixed finite jets differ from their limiting values by at most
\(C_jN^{C_j}(\eps+\lambda)\). In the intermediate window this
tends to zero uniformly. In particular
\(\nabla^2 G\le-(c_R/2)I\) for \(G=\Re\sum\phi_i\).

Write \(f_i=\partial_{\theta_i}\log Y\), \(g_i=\partial_{\theta_i}\sum\phi_i\),
and \(Q_s=\partial_s\sum\phi_i\). Permutation symmetry makes
\(D_{pq}=f_pg_q-f_qg_p\) zero at \(\theta_p=\theta_q\).
Division by the difference is justified for these smooth angular
functions by the integral of their transverse derivative; no analytic
extension of the deformation bumps is assumed.
At the limit it equals \(\ii[g'(\theta_q)-g'(\theta_p)]\).
The strict curvature and the finite-jet perturbation bound yield
\[
 D_{pq}=(\theta_p-\theta_q)C_{pq},\qquad |C_{pq}|\ge c_R/2 .
\]
The field
\(V_p=-Q_sf_q/D_{pq},\ V_q=Q_sf_p/D_{pq}\)
freezes \(Y,Z\). Its simple difference pole, divergence, and
rational real-angle weights have the two-power-per-Lie-order
bound of Appendix~\ref{app:lie}.

For separation \(d\ge N^{-K_0}\), all selected losses are polynomial
in \(N\), and the unchanged pair factors retain quadratic suppression.
For \(d<N^{-K_0}\), use \(\theta_i=t+\rho\omega_i\), with
\(\sum\omega_i=0,\ \sum\omega_i^2=1\).
Permutation symmetry gives \(\partial_\rho G(t,0,\omega)=0\),
so \(-\partial_\rho G\ge c_R\rho/2\).
The \(Y\)-phase is exactly \(Nt\); its real radial shifts do not
alter that phase. The coarea Jacobian in \(t,\rho\) is at least
\(c_RN\rho/2\). Convexity of the interval cube gives one interval
on each shape ray, and monotonicity gives multiplicity one there.

At a collision \(\theta_i=\theta_j\) the symmetric coupled shifts
are equal, so \(y_i-y_j\) vanishes to first order. The raw numerator
is bounded by \(C^{N^2}\rho^{N(N-1)}\), including its regular pair
denominators. After \(j\) Lie jets and one conservative flux/cutoff
loss, its remaining coarea power is
\(L=N^2-2j-4\ge4N_0+2>0\).
The exact inequalities \(d\le2\rho\) and \(\rho\le\sqrt N\,d\)
show that the near region
has \(\rho\le\sqrt N\,N^{-K_0}\), not literally \(\rho\le N^{-K_0}\).
Taking fixed \(K_0>1\) and increasing it if needed gives
\[
 C^{N^2}(\sqrt N\,N^{-K_0})^L
      e^{O_j(N\log(N+1))}(H+\log N)^2
 \le e^{C_jN\log(N+1)}(H+N)^{C_j}.
\]
The simple phase kernels have attenuation floors \(cN\eps\) in
this window. On each fixed shape ray, \(F=Nt\) and bounded individual
phases put the unwrapped image inside a rectangle with side lengths
at most \(CN\). The multiplicity-one statement above, followed by
the interval estimate in Appendix~\ref{app:lie}, bounds its phase
integral by \(CN^2(H+\log N)^2\). This explicit polynomial factor
is included in \(e^{O_j(N\log(N+1))}\) in the displayed estimate.
Counting the \(O(N^2)\) translated period boxes gives the same
factor; it is not an additional multiplicity to charge twice.
On the separated region the coarea \(1/\rho\) costs at most a
polynomial.
Selected-diagonal fluxes vanish by the rational numerator, and
full-equality fluxes vanish by the positive collision exponent at
fixed \(\eps\). The actual Stokes multiplier is smooth with bounded
finite jets and its \(q\) remains on a plateau. This proves the
claim uniformly in \(\lambda\le\lambda_*\).
\end{proof}

% PROV-INTERNAL: LARGE_LAMBDA_COMPLEX_DISK_AUDIT.md;
% PAIR_CONSTANTS_COMPLEX_DISK_FRESH_AUDIT.md.
\begin{lemma}[Protected large-current disks]\label{lem:protected}
On a named Stokes support, let \(\lambda_*\le\lambda\le1\).
For a sufficiently small constant \(c>0\), depending only on
the fixed point and previously chosen margins, the density is
holomorphic on
\[
 |s-s_\eps|\le\delta_s,\qquad \delta_s=c\lambda_*/N^2 .
\]
On this disk its product gaps are at least \(c'\lambda\),
its complete pairs have bound \(C^Ne^{-\kappa' N^2}\), and its
one-body factors have integrable majorants with total cost \(C^N\).
\end{lemma}
\begin{proof}
There is a named upper compact \(q\) with \(p_q=1\), so
\(P\ge1\). Fix a lower branch neighborhood inside \(m=1,p=0\)
and put \(a_B=\min|\sin\theta|>0\) there. With
\(t=\lambda P/N\), the exact imaginary-part formula gives
\[
 \Im W(s_\eps)\ge b_0\eps+\tfrac12\tau a_Bt
 \ge b_\tau(\eps+t),\qquad
 b_\tau=\min(b_0,\tau a_B/2)>0.
\]
Indeed the increase from the outward shift is at least
\(\tau a_Bt/2\), since \(\cosh v\ge1\) for real \(v\).
Write the disk constant as \(c_P\). After fixing \(\tau\), require
\[
 L_Sc_P\le b_\tau/4.
\]
Since \(t\ge\lambda_*/N\), throughout the disk
\[
 |\Delta W|\le L_Sc_P\lambda_*/N^2
 \le b_\tau t/(4N)\le b_\tau(\eps+t)/4.
\]
Thus these branch-near roots remain on the lower interior sheet.

Every non-plateau support angle has a fixed limiting distance
\(d_1>0\) from \(W=1\). Compute it on the preliminary endpoint
neighborhoods and the complement of the fixed lower plateau,
before shrinking \(\delta,\tau\); the upper branch is excluded.
All limiting
roots also have distance at least \(S_*\) from \(W=-1\).
For \(d=\min(d_1,S_*)>0\), the choices of
Lemma~\ref{lem:contraction}, followed by a reduction of \(c_P\),
ensure
\[
 \operatorname{dist}(W,\{1,-1\})
 \ge d-C(\eps+\tau+\delta)-L_S\delta_s\ge d/2.
\]
Here \(\delta\) is the fixed endpoint width.
Use the compact analytic continuations specified in that lemma.
Their motion is \(O(\delta_s)\); their modulus can exceed one.
The branch-coordinate motion is at most
\[
 C_\tau\frac{\delta_s}{\sqrt{\eps+t}}
 \le C_\tau c_P\frac{\sqrt{\lambda_*}}{N^{3/2}}.
\]
Reduce \(c_P\) further so these motions consume less than
the reserved \(q'-q_1\) and \(\eta/2\) pair margins and retain
the compact denominator margins. The regular double-branch
expression controls the remaining neighborhood. This proves
same-group bound \(q'\), cross bound \(1+\eta\), and the named
bound on the whole disk without a unit-disk Schur argument there.
The direct count in Lemma~\ref{lem:contraction} gives the same
quadratic exponent, including deletion of at most \(j\) pairs.

At the radial center the named upper root has
\(|z_q|\le e^{-c_q\lambda}\), \(c_q>0\), because its inward
shift is \(2\tau\lambda\) and its arc has positive sine.
Branch roots remain interior on the disk and compact roots obey
\(|z_i(s)|\le |z_i(s_\eps)|e^{C\delta_s}\).
Adding \(Cc_P\le c_q/2\) to the disk restrictions gives
\[
 |Z|\le e^{-c_q\lambda+CN\delta_s}\le e^{-c_q\lambda/2}.
\]
Also \(|Y|\le e^{-3\tau\lambda P/2}\), independently of \(s\).
Both product gaps are therefore at least \(c'\lambda\).
All restrictions on \(c_P\) are imposed after \(\tau\) and
before the final \(\eps_0\); an oversized disk need not preserve
the branch sheet. No current integral estimate has been used.

Finally put \(t=\lambda P/N\). The real displacement of the
branch center is \(O(\eps^2+t^2+\delta_s)\), and the imaginary
margin is at least \(c_2(\eps+t)\).
Split into \(|u|\le C(\eps+t)\) and its complement.
The imaginary margin bounds the square root in the first
case, and the linear real term bounds it in the second, with
the earlier choices of \(\tau\) and then \(c_P\) sufficiently small. Hence
\[
 |R(y;s)|\le C(|u|+\eps)^{-1/2}.
\]
This is a majorant independent of the coupled scalar \(P\);
no integration treats \(P\) as a free constant. Compact
one-body factors are bounded. The inverse-product numerator,
the exact Jacobian, and its Stokes replacement minor cost
\(C^N\). These facts prove all the stated assertions.
\end{proof}

\begin{proposition}[Large-\(\lambda\) current]\label{prop:largeS}
For \(N<C_H(H+1)^2\), the differentiated current integrated over
\(\lambda_*\le\lambda\le1\) satisfies \eqref{eq:reqLS}.
\end{proposition}
\begin{proof}
Lemma~\ref{lem:protected} bounds the absolute integral of
the undifferentiated density on its common disk by
\(C^Ne^{-\kappa N^2}\lambda^{-2}\), including the named-index
count in \(C^N\). The majorant is integrable and uniform on
that disk, so the integrated density is holomorphic there.
Cauchy's inequality of order \(j\) adds
\(C_jN^{2j}\lambda_*^{-j}\). Integration in \(\lambda\) costs
at most \(\lambda_*^{-1}\). Summing the convergent series
\(\sum_N C^NN^{2j}e^{-\kappa N^2}\) gives even the stronger
bound \(C_j\lambda_*^{-j-1}\). Its weaker version
\(C_j\lambda_*^{-j-2}\) is \eqref{eq:reqLS}.
By \eqref{eq:alphachoice}, this is \(o(\eps^{-1/2})\).
\end{proof}

% PROV-INTERNAL: SECOND_FAMILY_TAIL_GLOBAL_JET_AUDIT.md;
% SMALL_LAMBDA_LIE_ATLAS_AUDIT.md; PAIR_CONSTANTS_COMPLEX_DISK_FRESH_AUDIT.md.
\begin{proposition}[The remaining high orders]\label{prop:highKS}
The bound \eqref{eq:reqhigh} holds on
\(D\sqrt H\le N<C_H(H+1)^2\).
\end{proposition}
\begin{proof}
Use Lemma~\ref{lem:originaldisk}, which establishes the original
\(c\eps\) disk, the \(c\eps\) product gaps, integrable one-body
majorants, and complete pair bound \(C^Ne^{-\kappa N^2}\).
This invocation is independent of the large-current disk.
Cauchy of order \(j\) therefore gives
\(C_j^N\eps^{-j-2}e^{-\kappa N^2}\).
Integrating the small-\(\lambda\) current only reduces this bound.
The argument needs no enlarged Cauchy disk and makes no
unbounded-\(N\) assertion that \(1-r^N\ge cN\eps\).
\end{proof}

% END SOURCE


% BEGIN SOURCE: sections/09_natural_boundary.tex
\section{Summation and the natural-boundary implication}
\label{sec:conclusion}
% PROV-INTERNAL: SECOND_FAMILY_TAIL_GLOBAL_JET_AUDIT.md;
% SECOND_FAMILY_BULK_NATURAL_BOUNDARY_CANDIDATE.md;
% TAIL_COMPOSITION_HOSTILE_AUDIT.md (summation only; old central connector excluded).
\begin{theorem}[Uniform differentiated tail]\label{thm:tail}
For each \(0\le j\le k\),
\[
 \sum_{\substack{N>N_0\\N\ {\rm even}}}|T_N^{(j)}(\se)|
                =o(\eps^{-1/2}).
\]
\end{theorem}
\begin{proof}
The original full-order series is differentiated term by term at each
fixed exterior point; Appendix~\ref{app:pfaffian} supplies local normal
convergence of every fixed derivative. Only after differentiation are
the orders split at the \(\eps\)-dependent thresholds.
The ultra-high tail is negligible by Proposition~\ref{prop:ultrahigh}.
For the remaining orders use the exact identity \eqref{eq:FKS}.
The \(F\)-sum is \(\exp(O_j(\log^2H))\).
The large-\(\lambda\) current is at most
\(C_j\exp(\alpha_0(j+2)H)=o(e^{H/2})\).

Take a common positive \(\kappa\) smaller than all pair-contraction
constants, and then choose \(D\) with \(\kappa D^2>k+2\).
The high-\(N\) original and small-\(\lambda\) pieces have sum at most
\[
 \exp\bigl((j+2-\kappa D^2)H+O_j(\sqrt H)\bigr)=o(e^{H/2}).
\]
All intermediate pieces are covered by the fixed partition,
Propositions~\ref{prop:micro}, \ref{prop:allB}, \ref{prop:mixed},
and \ref{prop:compactR}. Their sum has logarithm at most
\(O_j(\sqrt H\log(H+2))=o(H)\).
The small-\(\lambda\) integration contributes at most its length
\(\lambda_*\le1\). The same choices work for all \(j\le k\), since
this is a finite set. This proves the claimed absolute tail estimate.
\end{proof}

% PROV-INTERNAL: SECOND_FAMILY_BULK_NATURAL_BOUNDARY_CANDIDATE.md, lines 31-41;
% FINAL_KILL_AUDIT_v1.md, lines 7-13.
% PROV-PRIMARY: TW2014 fixed-order smoothness and bulk expansion.
% RECONSTRUCTION: an explicit conditional implication, separate from sector validity.
\begin{theorem}[Bulk noncancellation criterion]\label{thm:conditional}
Assume the first-form-factor result of Theorem~\ref{thm:first} and
the absolute tail estimate of Theorem~\ref{thm:tail} at every
selected point. Then \(\T\) is a natural boundary of the exterior
bulk susceptibility germ.
\end{theorem}
\begin{proof}
Write \(U=\sum_{N\ {\rm even}}T_N\).
The finitely many terms below \(N_0\) have bounded \(k\)-th derivatives.
The first term has
\(T_{N_0}^{(k)}=2L_\beta\eps^{-1/2}+o(\eps^{-1/2})\).
The absolute higher-tail estimate excludes cancellation, giving
\[
 U^{(k)}(\se)=2L_\beta\eps^{-1/2}+o(\eps^{-1/2}).
\]
For \(j<k\), the finite first and lower terms are bounded and the
higher tails are \(o(\eps^{-1/2})\). Since \(\M^2\) is analytic and
nonzero near \(\sstar\), Leibniz's rule in \eqref{eq:bulk} therefore gives
\[
 \mathcal X^{(k)}(\se)
 =4\M^2(\sstar)L_\beta\eps^{-1/2}+o(\eps^{-1/2}),
 \qquad \M^2(\sstar)L_\beta\ne0.
\]
Holomorphic extension to a neighborhood of \(\sstar\) would bound this
derivative on a sufficiently short radial segment, a contradiction.
The selected points are dense in the upper-right arc; evenness and
conjugation give density in all four quadrants. A neighborhood of
any point of \(\T\), including an axis point, contains such a selected
point. Hence no neighborhood permits an exterior holomorphic extension.
The locally nonzero analytic factor \(\beta_{\mathrm{phys}}(s)\) transfers the same
conclusion to the physical susceptibility on its chosen continuation.
\end{proof}

% PROV-INTERNAL: SECOND_FAMILY_BULK_NATURAL_BOUNDARY_CANDIDATE.md;
% all active dependencies are itemized in the external source concordance.
% RECONSTRUCTION: consequence of the first coefficient and the uniform tail.
\begin{corollary}\label{thm:nb}
For the zero-field isotropic infinite square-lattice Ising model, the
unit circle in \(s=\sinh(2\beta_{\mathrm{phys}} J)\) is a natural boundary of the bulk
susceptibility germ continued from the low-temperature side.
\end{corollary}
\begin{proof}
Apply Theorem~\ref{thm:conditional} using Theorems~\ref{thm:first}
and \ref{thm:tail}.
\end{proof}

% END SOURCE

\appendix

% BEGIN SOURCE: appendices/A_cyclotomic_details.tex
\section{Cyclotomic field intersection}
% PROV-INTERNAL: TWO_NONZERO_PRIME_FAMILY.md, lines 16-40.
% RECONSTRUCTION: explicit special case, including the classical irreducibility argument.
\begin{lemma}\label{lem:cyclotomic}
If \(p\) is prime and \(p\nmid n\), then
\(\Q(\zeta_p)\cap\Q(\zeta_n)=\Q\).
\end{lemma}
\begin{proof}
We recall the needed elementary irreducibility argument.
The cyclotomic polynomial \(\Phi_m\) is monic with integer
coefficients, obtained successively from \(X^m-1=\prod_{d\mid m}\Phi_d\).
Let \(f\) be the monic irreducible polynomial of a primitive \(m\)-th
root \(\zeta\). By Gauss's lemma, \(f\in\Z[X]\) divides \(\Phi_m\).
For a prime \(q\nmid m\), suppose \(\zeta^q\) belonged to a different
irreducible factor \(g\) of \(\Phi_m\). Then \(f\mid g(X^q)\).
Modulo \(q\), \(g(X^q)=g(X)^q\); hence the reductions of \(f\)
and \(g\) share a factor. Their product divides \(X^m-1\) modulo
\(q\), making that polynomial have a repeated factor.
But its derivative is \(mX^{m-1}\), coprime to \(X^m-1\)
in characteristic \(q\), a contradiction.
Thus taking a \(q\)-th power preserves the conjugacy class for
every prime \(q\nmid m\). Factoring any positive integer coprime
to \(m\) into such primes shows that every primitive \(m\)-th
root is a conjugate. Therefore \(\Phi_m\) is irreducible and
\([\Q(\zeta_m):\Q]=\varphi(m)\).

For \(p\nmid n\), the compositum
\(\Q(\zeta_p,\zeta_n)=\Q(\zeta_{pn})\), by Bezout's identity.
Its degree is
\(\varphi(pn)=\varphi(p)\varphi(n)\).
This is the maximal possible product degree. If the intersection
contained a field of degree greater than one, a basis calculation
over that intersection would make the compositum degree strictly
smaller than this product. The intersection is therefore \(\Q\).
\end{proof}

For the sparse relation used in Theorem~\ref{thm:prime}, a symmetric
polynomial representing \(c_0+\sum c_j(\zeta^j+\zeta^{-j})\) has
degree at most \(p-1\). Its vanishing makes it a scalar multiple
of \(\Phi_p\). Comparing its paired coefficients gives precisely
\(c_0=c_1=\cdots=c_d\). The requirement \(d\ge5\) is what rules
out a nonzero relation involving at most four nonconstant positions.

% END SOURCE


% BEGIN SOURCE: appendices/B_residue_and_local_model.tex
\section{The order of local integrations}\label{app:local}
% PROV-INTERNAL: COMPACT_PRIME_MEAN_RESIDUE_AMPLITUDE.md, lines 31-69.
In the homogeneous first-chart scaling \(v=\eps w\),
\(t=\sqrt\eps\,\tau\), the two global denominators are
\[
 \eps(Nc_0-\ii w),\qquad
 \eps(NA_0+\ii bw-\ii dR^2).
\]
The first pole is below and the second above the real \(w\)-line.
For fixed \(R^2=|\tau|^2\),
\[
 \int_\R
 \frac{\dd w}{(Nc_0-\ii w)(NA_0+\ii bw-\ii dR^2)^{k+1}}
 =\frac{2\pi}{(N(A_0+bc_0)-\ii dR^2)^{k+1}}.
\]
The combined damping is \(N(A_0+bc_0)=Q\), independent of the arbitrary
evaluation contour radius. The shape volume scales as
\(\eps^{(N-1)/2}\), the Vandermonde as
\(\eps^{N(N-1)/2}\), and the mean measure as \(\eps\).
The two denominator powers are \(\eps^{k+2}\).
The net exponent is \(-1/2\).

The absolute joint integral before mean integration is not a permissible
dominated-convergence majorant: its high-order \(Z\)-ridge moves with
\(R^2\). The manuscript obtains an exact mean residue on a fixed physical
rectangle first. Only its constrained shape integral has the integrable
\(R^{-2}\) tail used for the limit.

For completeness, \(\mathscr C_N>0\) needs no tabulated Mehta formula.
The squared Vandermonde is nonnegative on the real shape sphere and
positive in a neighborhood of any zero-sum vector with distinct
coordinates. That sphere has finite measure. Thus the beta period's
nonvanishing follows from an ordinary positive angular integral and
nonzero gamma factors. No numerical coefficient is used.

% END SOURCE


% BEGIN SOURCE: appendices/C_pfaffian_bounds.tex
\section{Pfaffian identities and exterior convergence}\label{app:pfaffian}
% PROV-INTERNAL: PFAFFIAN_TAIL.md; L2_PAIR_TAIL.md; RESIDUE_REDUCTION.md.
% RECONSTRUCTION: rational-function induction rather than finite numeric verification.
The following classical Schur identity and Hadamard's determinant
inequality are standard tools, not originality claims. The needed
Schur identity is proved here to fix its sign and convention.
\begin{lemma}[Schur identity]\label{lem:schur}
For even \(N\) and \(A_{ij}=(v_i-v_j)/(1-v_iv_j)\),
\[
 \Pf A=\prod_{i<j}\frac{v_i-v_j}{1-v_iv_j}.
\]
\end{lemma}
\begin{proof}
Set \(x_i=(1+v_i)/(1-v_i)\), so the matrix entries are
\((x_i-x_j)/(x_i+x_j)\).
Induct on even \(N\), regarding both sides as rational functions of
\(x_N\). At a simple pole \(x_N=-x_i\), Pfaffian expansion in the
last index gives residue
\((-1)^{i+1}2x_i\) times the \((N-2)\)-Pfaffian.
On the product side, each index between \(i\) and \(N\) contributes
one minus sign when its two paired factors are evaluated there;
the total sign is \((-1)^{N-i-1}=(-1)^{i+1}\).
The residues agree by induction. The difference has no finite poles
and is bounded at infinity, so is constant. At \(x_N=x_{N-1}\),
both expressions vanish, fixing that constant to zero.
Degenerate cases follow as rational identities by continuation.
\end{proof}

If an even \(m\)-by-\(m\) antisymmetric matrix has entries at most \(C\),
then Hadamard's determinant inequality and
\((\Pf A)^2=\det A\) give
\(|\Pf A|\le C^{m/2}(m-1)^{m/4}\).
The number of perfect matchings is
\((N-1)!!=N!/(2^{N/2}(N/2)!)\).
These are different bounds; the selected-contour proof uses bounded
minors and an explicit expansion only for exceptional compact pairs.

Fix a compact exterior neighborhood \(K\) and a slightly larger
compact exterior neighborhood \(K_+\) with
\(K\subset\operatorname{int}K_+\). Choose an admissible radius
\(r<1\) and \(q<r\) such that every interior root satisfies
\(|z(y;s)|\le q\) for \(s\in K_+\) and \(|y|=r\).
On this compact set the dispersion quantity is bounded, say
\(|W(y;s)|\le M_W\). The equation \(z+z^{-1}=2W\) gives
\[
 |z|^{-1}=|2W-z|\le 2M_W+q,
 \qquad m:=(2M_W+q)^{-1}\le |z|.
\]
Thus individual one-body factors and pair matrix entries have bounds
independent of \(N\), but the global inverse products need an
exponential bound. In fact, \(|Y^{-1}|=r^{-N}\), and
\[
 |Z^{-1}+Y^{-1}|\le m^{-N}+r^{-N}\le 2C_1^N,
 \qquad C_1:=\max\{m^{-1},r^{-1}\}.
\]
For even \(N\ge2\), the global gaps satisfy
\(|1-Z|\ge1-q^N\ge1-q^2\) and
\(|1-Y|\ge1-r^N\ge1-r^2\).
Applying Hadamard's inequality to the two Pfaffians in
\eqref{eq:reduced}, and absorbing the bounded one-body factors and
contour measures into a constant to the power \(N\), yields
\[
 |T_N(s)|\le C_0\frac{C_2^N N^{N/2}}{N!}
       \le C_0\left(\frac{eC_2}{\sqrt N}\right)^N,
 \qquad s\in K_+.
\]
Here \(C_0,C_2\) are independent of \(N\), and the last inequality
uses \(N!\ge(N/e)^N\).
Choose \(d>0\) such that each closed \(d\)-disk centered in \(K\)
lies in \(K_+\). Cauchy's inequality gives, for every fixed
nonnegative integer \(j\),
\[
 \sup_{s\in K}|T_N^{(j)}(s)|
 \le j!d^{-j}C_0\left(\frac{eC_2}{\sqrt N}\right)^N.
\]
The majorant is summable over even \(N\), proving local normal
convergence of every fixed differentiated series and justifying
termwise differentiation in the exterior. These compact-set
constants need not stay bounded on approach to the unit circle;
this argument does not supply the uniform boundary-tail estimates.

% END SOURCE


% BEGIN SOURCE: appendices/D_lie_fields_and_coarea.tex
\section{Fixed-cutoff differentiation, pole counting, and coarea}
\label{app:lie}
% PROV-INTERNAL: LIE_POLE_FILTRATION_PROOF.md;
% ORIGINAL_CONTOUR_TWO_PHASE_FIELD.md; ALL_B_LIE_NUMERATOR_DESOURCE.md;
% SECOND_FAMILY_TAIL_GLOBAL_JET_AUDIT.md.
% V2 EXPANSION: unfactored jet recurrence, chartwise identity, and flux order.

\subsection{The identity on a fixed real integration domain}
\begin{lemma}[Fixed-domain differentiation]\label{lem:lie}
Let \(w\) be a smooth, \(s\)-independent weight on a real angular torus.
For an \(s\)-dependent complex vector field \(V\) and a top form \(A_s\),
set \(\mathcal D=\partial_s-\Li_V\). If all puncture fluxes vanish
at every intermediate stage through order \(j\), then
\[
 \partial_s^j\int wA_s=\int\mathcal D^j(wA_s).
\]
If \((\partial_s-V\cdot\nabla)Y
=(\partial_s-V\cdot\nabla)Z=0\), the factors
\((1-Y)^{-1}\) and \((1-Z)^{-1}\) remain simple throughout this identity.
\end{lemma}
\begin{proof}
On a top form, \(\Li_V=\dd\iota_V\) is a divergence. Its integral is the
boundary flux, including every puncture boundary. Its vanishing proves
the first derivative identity. Apply that identity to the entire
\(s\)-dependent top form produced at the preceding step. In particular,
the iteration differentiates \(V\), its divergence, and all coefficients
already present. The kernel assertion follows from the derivation rule.
Complex coefficients in \(V\) are allowed by linearity of real
integration by parts. No real flow or holomorphic extension of \(w\)
is required.
\end{proof}

If \(\sum_h w_h=1\) is a fixed smooth partition, use a different field
\(V_h\) on each weighted term:
\begin{equation}
 \partial_s^j\int A_s
   =\sum_h\int(\partial_s-\Li_{V_h})^j(w_hA_s).
 \label{eq:chartwiseLie}
\end{equation}
No equality between the fields on overlapping supports is required.
The proof is the preceding lemma applied to each integral, followed
by addition. This observation avoids interpolating fields in a way
that could destroy either frozen-product identity.

\subsection{Regular branch coordinates and the selected-pair field}
On the original all-branch chart let \(\chi_s:u\mapsto\phi\), where
\(\phi=\arccos W\). The inverse \(y_s(\phi)\) is even and analytic
on a fixed polydisk about \((s_*,0)\), because \(W_\theta\ne0\) there.
The one-body form and the complete pair are
\begin{equation}
 \frac{\dd y}{2\pi\ii}\frac{z}{\sinh\gamma}
   =-\frac{z}{\pi(1-y^{-2})}\dd\phi,\qquad
 P_{il}=(\phi_i-\phi_l)^2 A_s(\phi_i,\phi_l).
 \label{eq:branchregular}
\end{equation}
Here \(A_s\) is analytic with bounded fixed jets. Indeed,
\(y_s(\phi_i)-y_s(\phi_l)\) has the two factors
\(\phi_i-\phi_l\) and \(\phi_i+\phi_l\); the latter cancels the
corresponding factor in \(1-e^{-\ii(\phi_i+\phi_l)}\).
No division by a remaining collision zero is used.

Write the regular top form as
\[
 \Omega_s=
  \frac{\mathscr N_s(\phi)}{\mathcal K_s(\phi)}
       \dd\phi_1\cdots\dd\phi_N,\qquad
 \mathcal K_s=(1-Y_s)(1-Z),\qquad
 \mathscr N_s=\Delta(\phi)^2\mathscr A_s(\phi).
\]
The analytic factor \(\mathscr A_s\) includes all normalized measures,
the inverse-product numerator, and the regular pair coefficients.

Put \(a_i=S'/W_{u_i}\), \(b_i=\partial_{u_i}\phi_i\), and
\(A=\sum_i a_i\). For a selected pair \(p,q\), set
\begin{equation}
 \begin{split}
 V_i&=a_i\quad(i\ne p,q),\\
 V_p&=a_p+\frac{Ab_q}{b_p-b_q},\qquad
 V_q=a_q-\frac{Ab_p}{b_p-b_q}.
 \end{split}\label{eq:twofield}
\end{equation}
Then \(\sum_iV_i=0\) and
\(\sum_i b_iV_i=\partial_s\sum_i\phi_i\), so \(Y,Z\) are frozen.
In regular coordinates the residual field
\(B=\partial_s\chi_s-(\chi_s)_*V\) has only two nonzero entries:
\[
 B_p=-\frac{Ab_pb_q}{b_p-b_q},\qquad B_q=-B_p.
\]
Since \(b(\phi)=B_s(\phi)/\phi\) with \(B_s\) even, analytic, and
nonzero near zero,
\[
 \frac{b_pb_q}{b_p-b_q}
     =\frac{H_s(\phi_p,\phi_q)}{\phi_q-\phi_p}.
\]
The analytic coefficients and any fixed number of their derivatives
have polynomial \(N\)-bounds. Both \(V\), expressed in regular
coordinates, and \(B\) have at most one selected-difference pole;
\(\operatorname{div}_\phi B\) has at most two.

The pullback identity needed for smooth real cutoffs is
\begin{equation}
 (\partial_s-\Li_V)(w\chi_s^*\Omega_s)
  =w\chi_s^*(\partial_s+\Li_B)\Omega_s
       -(V\cdot\nabla_u w)\chi_s^*\Omega_s.
 \label{eq:pullbackLie}
\end{equation}
It follows by differentiating the pullback and subtracting the
pushforward of \(V\). In particular,
\((\partial_s+B\cdot\nabla_\phi)\mathcal K_s=0\).
Formula \eqref{eq:pullbackLie} is used on the real \(u\)-domain;
the smooth function \(w\) is never extended holomorphically in \(\phi\).

\subsection{The complete high-order recurrence}
\begin{lemma}[Unfactored jet filtration]\label{lem:jets}
For each fixed \(j\), the \(j\)-fold operation in
\eqref{eq:pullbackLie} is a finite sum of terms
\begin{equation}
 (\partial_u^\mu w)(u)\,
 \chi_s^*\left[
  \frac{c_{j,\mu,\ell,\nu}(s,\phi)\,
        \partial_s^\ell\partial_\phi^\nu\mathscr N_s(\phi)}
       {\mathcal K_s(\phi)}
       \dd\phi_1\cdots\dd\phi_N\right].
 \label{eq:jetexpansion}
\end{equation}
If \(m\) is the selected-difference pole order of a displayed
coefficient, then that term can be chosen to satisfy
\begin{equation}
 m+|\mu|+|\nu|\le2j,\qquad
 \ell+|\mu|+|\nu|\le j .
 \label{eq:jetdegree}
\end{equation}
Apart from these explicitly counted poles, fixed coefficient jets
are bounded by \(C_jN^{C_j}\) on the fixed regular-coordinate
neighborhood. There are at most \(C_jN^{C_j}\) terms.
\end{lemma}
\begin{proof}
For \(j=0\) take \(c=1\). Suppose \eqref{eq:jetexpansion} holds.
Apply \eqref{eq:pullbackLie} to each term, and use the fact that
the kernel is frozen. The possibilities are as follows:
\begin{enumerate}
\item A fixed-\(\phi\) \(s\)-derivative on \(c\) does not raise its
difference-pole order: \(\delta=\phi_p-\phi_q\) is independent of
\(s\) in these coordinates.
\item \(B\cdot\nabla_\phi c\), or multiplication by
\(\operatorname{div}_\phi B\), raises \(m\) by at most two.
\item \(B\cdot\nabla_\phi\) on the numerator derivative raises
\(|\nu|\) by one and \(m\) by at most one.
\item \(\partial_s\) on the numerator derivative raises \(\ell\)
by one without adding a difference pole.
\item The final term of \eqref{eq:pullbackLie} raises \(|\mu|\)
by one and contributes one component of \(V\), whose pole
order is at most one.
\end{enumerate}
Every operation therefore raises the first count by at most two
and the second count by at most one. These are all product-rule
terms, including derivatives of the previous field coefficients.
Bounded analytic coefficients, at most \(N\) summands in \(A\),
and a fixed number \(j\) of differentiations give the claimed
polynomial term and coefficient bounds.
\end{proof}

The object differentiated in this recurrence is the
\emph{unfactored numerator} \(\mathscr N_s\).
A numerator derivative of order \(\ell+|\nu|\le j\) changes at
most \(j\) pair factors. The remaining pair contractions survive.
Fixed-\(\phi\) \(s\)-derivatives preserve the Vandermonde polynomial;
\(|\nu|\) coordinate derivatives reduce its full-collision degree
by at most \(|\nu|\).
Dividing the answer by the original Vandermonde would instead
create poles at unselected collisions and would invalidate the
selected-pole statement. Such division is not part of
Lemma~\ref{lem:jets}.

\subsection{Rational weights and all boundary fluxes}
Fix an integer \(M>2k+2\), independent of \(N\), and put
\[
 d=\max_{i<l}|u_i-u_l|,\qquad
 w_{pq}=\frac{(u_p-u_q)^{2M}}
                  {\sum_{i<l}(u_i-u_l)^{2M}}\quad(d>0).
\]
They sum to one off full equality. For \(|\mu|\le k\),
\begin{equation}
 |\partial_u^\mu w_{pq}|
 \le C_kN^{C_k}
       \frac{|u_p-u_q|^{2M-|\mu|}}{d^{2M}}.
 \label{eq:weightjets}
\end{equation}
To see this, differentiate the quotient. A derivative of the
denominator contributes a power at most \(d^{-1}\); replace it,
when needed, by \(|u_p-u_q|^{-1}\). The denominator itself is
at least \(d^{2M}\), and its derivative sums have polynomial
\(N\)-cost. This proves \eqref{eq:weightjets} term by term.

On the original branch cone,
\(|\phi_p-\phi_q|\ge c|u_p-u_q|\), by
Lemma~\ref{lem:branch}. Combining \eqref{eq:weightjets} with
\eqref{eq:jetdegree} absorbs every selected pole at a cost
\(C_jN^{C_j}d^{-2j}\) on a separated support, while the
undifferentiated pair factors retain their quadratic contraction.
Other fixed smooth cutoffs are included by the ordinary product
rule. A cutoff at scale \(b_N\) or \(\rho_N\) has its corresponding
inverse-scale cost; these costs are retained in
Propositions~\ref{prop:micro} and \ref{prop:allB}.

Near full equality, the numerator has degree \(N(N-1)\).
The bound \(m+|\mu|+|\nu|\le2j\) shows that a full \(j\)-th
density loses at most \(2j\) powers of the shape radius.
A further inverse radius is reserved when a boundary flux or
a transition-cutoff estimate requires it. This is the origin
of the conservative loss \(2j+1\) used in the sector proofs.

These facts also verify the hypotheses of Lemma~\ref{lem:lie};
they are not additional terms to discard.
First take a selected-diagonal tube of radius \(h\) on a region
where \(d\) is bounded below. The pair double zero, the
\(2M\)-power weight, the preceding \(j\) operations, and one
extra field factor give a flux at most
\[
 C_{\eps,N,j}\,h^{2M+1-2j}.
\]
It vanishes. Next take a full-equality sphere of radius \(\rho\).
At fixed \(\eps>0\) all global denominators and regular coordinate
changes are bounded; the shape surface measure contributes
\(\rho^{N-2}\). A sufficient flux exponent is
\[
 N(N-1)-2j-1+N-2=N^2-2j-3>0
 \quad(N\ge N_0+2,\ j\le k).
\]
The limits are taken in this order at fixed \(\eps,N\).
There are no boundaries at a cutoff support endpoint because
the cutoffs are smooth on the closed torus. These flux checks
apply at every earlier stage \(0\le l\le j\), and establish the
iterated identity before any boundary-uniform estimate is taken.

\subsection{Mixed and compact charts}
In a mixed chart only the true branch coordinates are replaced
by \(\phi_i\). The selected denominator is \(1-b_q/b_j\), bounded
away from zero on the nested support, and \(1/b_j\) is analytic
at the branch. There is consequently no selected-collision pole
in this calculation. The same pullback recurrence includes the
regular branch one-body form, the compact-angle measures, and
the full contour Jacobian. All fixed jets have polynomial
\(N\)-cost, apart from \(C^N\) one-body products.

For the compact right chart the coefficients are smooth in real
angles and analytic in \(s\). The factorization
\(D_{pq}=(\theta_p-\theta_q)C_{pq}\), with \(|C_{pq}|\ge c>0\),
uses smooth division, as proved in Proposition~\ref{prop:compactR}.
The same five product-rule cases prove a two-power loss per
operation in the real angular variables. There are no singular
branch-coordinate changes in this compact case. It does not
require analytic continuation of a deformation bump.

\paragraph{Coupled occupancy and changed pairs (rc4 repair).}
The compact-right current field may move a coordinate where
\(p'\ne0\). The exact contour Jacobian gives
\[
 (V\cdot\nabla\log y_i)_{\rm occ}
   =\frac{\lambda\tau}{2N}m_i\sum_l V_l p'_l.
\]
An individual off-diagonal entry is \(O(\lambda\tau/N)\).
The compact-right field has two nonzero components, each possibly
\(O(N/|\theta_p-\theta_q|)\), so the contraction can be
\(O(\lambda\tau/|\theta_p-\theta_q|)\). It moves every
\(m\)-supported root and must be retained. All these roots are
compact in this chart; no singular raw branch-angle jet is used.

To count descendants, keep each complete pair \(F_e=P_{il}\)
as a scalar function of all angles and \(s\), including its
occupancy dependence. Write the density as \(G\prod_e F_e\).
The scalar operator \(\partial_s-V\cdot\nabla\) is a derivation;
the top-form correction is multiplication by
\(-\operatorname{div}V\). In one product-rule summand a step
hits a coefficient or \(G\), an already changed pair, or one
previously unchanged pair. Divergence changes none. Induction
therefore leaves at most \(j\) changed pairs per descendant at
order \(j\), even though the sum of descendants involves all pairs.
In regular coordinates, hitting either the double-zero polynomial
or its analytic coefficient marks that same pair. No division
by an original pair is involved.

There are \(O(N^2)\) possible pair choices at a step. Expanding a
coordinate contraction or an occupancy sum introduces at most
\(O(N)\) further choices. A fixed number of chain rules therefore
has \(N^{O(j)}\) descendants. Fixed jets of \(P/N\) are bounded
by the fixed bump jets, and fixed compact root and pair jets are
bounded on the continuation neighborhoods of
Lemma~\ref{lem:contraction}. Sums of phase jets, and inverses of
the separated smooth divided determinant \(C_{pq}\), have only
polynomial fixed-jet costs. The explicit selected difference pole
still adds at most two powers per Lie step. One-body and Jacobian
products cost \(C_j^N N^{O(j)}\). Consequently the coupled terms
fit the existing polynomial coefficient and term budgets and
the unchanged pairs retain \(C_j^N e^{-\kappa N^2}\).
In mixed Stokes charts the plateau identities
\(V\cdot\nabla P=V\cdot\nabla\sum m_i=0\) remove this coupling;
original \(K\) has \(\lambda=0\).

\subsection{Simple kernels and the order of coarea}
For \(a>0\),
\[
 \frac1{|1-e^{-a+\ii t}|}
 \le\frac C{\min(1,a)+\operatorname{dist}(t,2\pi\Z)}
 \quad\bigl(0\le\operatorname{dist}(t,2\pi\Z)\le\pi\bigr).
\]
On an interval of length \(L\), the integral is at most
\(C(1+L)(1+|\log a|)\) for \(0<a<1\). A variable damping can be
replaced by its positive lower bound in this estimate.

For two unwrapped phases \(F,G\), suppose their image is contained
in \(I_F\times I_G\), the map has multiplicity at most \(m_0\),
and its real Jacobian has already been divided out in the density
bound. If the attenuation floors are \(a_Y,a_Z\in(0,1)\), the
remaining two-kernel integral is bounded by
\[
 C m_0(1+|I_F|)(1+|I_G|)
       (1+|\log a_Y|)(1+|\log a_Z|).
\]
Indeed extend the nonnegative integrand to the containing rectangle
and apply the one-dimensional bound twice. Equivalently, subdivide
each interval into translates of one period. This accounts for
periodic wrapping separately from the multiplicity of the unwrapped
map. In the near-equality sectors \(m_0=1\) on each fixed shape
ray and both lengths are at most \(CN\), yielding an explicit
\(CN^2\) factor. Any already counted index cells or fixed-jet
terms retain their stated polynomial costs. Such factors are
absorbed by the existing \(e^{C_jN\log(N+1)}\) budgets.

Coarea is applied only after all differentiation, to maps whose
real Jacobian and multiplicity are specified in the sector proof.
Sign and extreme-index cells are introduced only in that final
absolute integral. They are not nonsmooth cutoffs in
\eqref{eq:chartwiseLie}. When a shape radius is exchanged for a
phase variable, its remaining positive power is bounded on the
support before integrating the phase. It is not integrated a
second time as though it were still an independent variable.

% END SOURCE


% BEGIN SOURCE: appendices/E_constant_dependencies.tex
\section{Constant and quantifier dependencies}\label{app:constants}
The constants can be selected in the following order. Each step
uses only quantities selected earlier, not the final value of
\(\eps\) or the particle number.
\begin{enumerate}
\item Fix \((p,a,b)\), hence \(N_0,k,S_*,\theta_B\), and the
original radial constant and a preliminary bound on \(|S'|\).
The final disk constants are chosen in step 4. There are only finitely many
derivative orders \(0\le j\le k\).
\item Fix a sufficiently small branch neighborhood and compact
margins. Choose the nested ratio \(L\) large enough for the
outer-anchor/true-branch slope separation, also relative to the
bounded upper-current slopes on a preliminary endpoint neighborhood
separated from the upper branch. Fix the smooth nested
partition and all of its finite derivative bounds.
\item Obtain the same-group pair constant \(q<1\) at the limiting
real support, with the groups and continuation neighborhoods of
Lemma~\ref{lem:contraction}. Choose \(q<q_1<q'<1\), then the cross-group slack
\(\eta\), and then the upper endpoint slivers and deformation
strength \(\tau\), retaining the reserved same/cross disk slack.
The upper branch must lie inside \(a=p=1\). Fix the bump jets
and the positive constants \(b_\tau,c_q\) at this step.
\item Choose the selected \(c/N\), original \(c\eps\), and
protected-current \(c\lambda_*/N^2\) disk constants small enough
for the already fixed margins. In particular impose
\(L_Sc_O\le b_0/4\), \(L_Sc_P\le b_\tau/4\), the
compact-distance and reserved pair-slack restrictions, and
\(Cc_P\le c_q/2\) from Lemma~\ref{lem:protected}.
No reduction of \(\tau\) follows this choice. Take a common smaller positive
\(\kappa\) for the finitely many pair estimates.
\item Choose \(C_H\) from the original Pfaffian bound,
\(0<\alpha_0<1/[2(k+2)]\), and then \(D\) with
\(\kappa D^2>k+2\). Fix the rational-weight integer \(M>2k+2\)
and a sufficiently large compact-equality exponent \(K_0>1\).
\item Choose the resultant exponent \(A\), followed by the small
microcore prefactor \(c_*\) in \(b_N=c_*e^{-(A+4)N}/N\).
Determine the finite-jet constants on the anchor neighborhood
\(\rho\le b_N/8\), \(|u_a|\ge b_N/2\). Its lower bound on
\(|t|\) determines the cubic coefficient \(C\) independently of
the still unselected exponent \(B\).
\item With the finite \(\Gamma_*\) defined in
Proposition~\ref{prop:allB}, choose
\(B,C_1>A+4+\max\{1,\Gamma_*\}\). This ensures the geometric
inclusions \(2\rho_N\le b_N/8\) and \(\delta_N\le b_N/16\)
for every \(N\ge N_0+2\), including the finitely many low orders.
Increase \(B\) further so that
\(\exp((A+4)N^3/2+\tfrac12N^2\log N+C_jN^2)
 (2e^{-BN})^{N^2-2j-4}\) is controlled for all large
\(N\) and every \(j\le k\). Only the remaining numerical bounds
at finitely many orders are absorbed into a point-dependent
constant; their support geometry already satisfies the inclusions.
\item Finally decrease a single \(\eps_0>0\) for all the conditions.
In the intermediate window, for every fixed \(C,m,r\),
\[
 \eps^r e^{CN}N^m\longrightarrow0\quad(r>0),\qquad
 \eps/b_N^m\longrightarrow0,
 \quad N\le D\sqrt{\log(1/\eps)}.
\]
These limits also give \(N^{C_j}(\eps+\lambda_*)\to0\).
They are used only in that window.
In particular impose \(N\eps e^{AN}\le c\),
\(\eps/b_N\le c\), and the finitely many compact-jet
perturbation inequalities for \(j\le k\), as well as
\(\eps+\lambda_*\) sufficiently small relative to the fixed
branch and compact margins. For the first two the logarithms
at \(N\le D\sqrt H\) are bounded by
\(-H+O(\sqrt H+\log H)\); for the compact perturbations the bound
is \(-\alpha_0H+O(\log H)\). Thus one \(\eps_0\), independent
of \(N\), enforces every required condition. Only finitely many
fixed exponents from the preceding estimates are needed.
\end{enumerate}

The distinction between costs is essential: an ordinary one-body
product costs \(C_j^N\); a full collision coefficient may cost
\(e^{C_jN^2}\); conversion from the exponentially small branch
anchor to shape coordinates may cost \(e^{CN^3}\). The last cost
is paid only where the positive collision power at \(\rho_N\)
absorbs it. It is not replaced by a polynomial in \(N\).

\begin{longtable}{>{\raggedright\arraybackslash}p{.22\textwidth}>{\raggedright\arraybackslash}p{.30\textwidth}>{\raggedright\arraybackslash}p{.38\textwidth}}
\toprule
Quantity & May depend on & Uniformity actually used\\
\midrule
\endhead
\(c_0,c_1\), original radius/disk &
Fixed selected point & All original angles and all even \(N\);
radius held fixed under each \(s\)-derivative.\\
Branch and compact margins &
\(S_*,\theta_B,k\) & Independent of \(N,\eps\); compact groups
exclude the exact branch before strict constants are chosen.\\
\(q,\eta,\kappa,\tau\) &
The preceding margins and point & Entire stated support, all
occupancies \(P/N\in[0,1]\); radial centers and both stated disks.\\
Selected-contour disk \(c/N\) &
Point, margins, \(\tau\) & All \(\eps\) sufficiently small;
not an assertion about the original undeformed contour on that disk.\\
Current disk \(c\lambda_*/N^2\) &
Point, margins, \(\tau\) & \(\lambda_*\le\lambda\le1\);
compact roots may leave the unit disk.\\
\(C_j\) in jets &
Point, fixed bumps, \(j\le k\) & Independent of \(N,\eps\) within
the stated window; powers \(N^{C_j}\) are explicit.\\
Final \(C_j\) in \eqref{eq:reqinter} &
Point, \(j\), all preceding margins, \(\tau,\kappa,\alpha_0,D\),
\(M,K_0,A,c_*,B,C_1\) & Independent of \(N,\eps,\lambda\).
Finite small-order constants may be large but are
\(\eps\)-independent.\\
\(A,b_N,\rho_N\) &
Algebraic point; \(B\) also on \(k\) & Exponential \(N\)-dependence
is retained. \(\eps e^{cN}\to0\) only for \(N\le D\sqrt H\).\\
\(\alpha_0,D,C_H\) &
Point and \(k\) & Fixed before \(\eps\to0\); all lower derivatives
use the same finite choices.\\
First-chart constants &
The fixed \(N_0,p,a,b\) & No uniform first-amplitude estimate
in \(p\) or \(N_0\) is claimed.\\
\bottomrule
\end{longtable}

Every estimate stated for sufficiently small \(\eps\) means
\[
 \forall(p,a,b)\ \exists\eps_0,C_j,\ldots\
 \forall\,0<\eps<\eps_0\
 \forall\,N\text{ in the specified window}.
\]
Neither the point nor the derivative order is placed after the
\(\eps\)-quantifier. The \(N\)-window and \(\lambda\)-window splits
are made after differentiation; their moving endpoints are not
being differentiated.

% END SOURCE


% BEGIN SOURCE: appendices/F_normalization_checks.tex
\section{Normalization and the independent angular representation}
\label{app:normalization}
% PROV-INTERNAL: FORM_FACTOR_NORMALIZATION_AUDIT.md;
% ONSITE_GROUPING_EXACT_AUDIT.md;
% COMPACT_FOUR_CHART_NORMALIZATION_RESOLUTION.md.
% PROV-PRIMARY: TW2014 Appendix 1 (7), its normalized y-residue identity;
% Boukraa et al. 2008 (4)-(9); McCoy and Maillard 2012 (67)-(70), comparison only.
The squared formula in Theorem~4.1 of \cite{PalmerTracy1981}
reduces on p.~375 to \(\det_2(I+G)=\det_2(I+g)^2\).
For the smooth trace-class two-point kernel, the zero diagonal gives
\(\operatorname{Tr}g=0\), hence \(\det_2(I+g)=\det(I+g)\).
The determinant is real and continuous for real \(s>1\);
the Griffiths bound used on p.~359 keeps its modulus at least one,
and its limit one at infinity fixes the positive square root.
The resulting unsquared series is Theorem~5.0 and its proof on
pp.~374--378. In the representation of the published Appendix~1
of \cite{TW2014}, it has
\(N\) factors \(\dd\theta/(2\pi)\).
With \(x=e^{\ii\theta}\), these become \(\dd x/(2\pi\ii x)\).
Introducing \(N\) auxiliary residues uses another \(N\) factors
\(\dd y/(2\pi\ii)\). Thus the double contour prefactor is
\((2\pi\ii)^{-2N}/N!\) for ordinary unnormalized contours,
as displayed with the published equation~(3) of \cite{TW2014}.
Summing the nonorigin lattice sites gives
\[
 \frac4{(1-X)(1-Y)}-\frac2{1-X}-\frac2{1-Y}
 =\frac{2(X+Y)}{(1-X)(1-Y)}.
\]
The factors \(1/(XY)\) in the normalized correlation measures
then give the numerator in \eqref{eq:double} and the factor two
in \eqref{eq:bulk}; its physical origin term is
\(\langle\sigma_{00}^2\rangle-\M^2=1-\M^2\).

For fixed real \(s>1\), choose a small common angle strip with
\(\Re\gamma\ge c>0\) and bounded canceled kernel. Shifting all angles
by \(\ii\eta\operatorname{sgn}(a)\) and applying Hadamard bounds the
\(N\)-particle coefficient at the lattice site \((a,b)\) by
\[
 \frac{N^{N/2}C^N}{N!}
 e^{-N(\eta|a|+c|b|)}.
\]
The determinant integrand is periodic, and the bound is summable in
sites and even \(N\ge2\). This justifies the lattice/Fredholm
interchange on the real interval; it supplies no boundary-uniform
infinite-tail estimate.

For fixed even \(N\), define \(C_N^{\rm std}\) as the full-site
coefficient \(\chi^{(N)}\) displayed with the published equation~(3)
of \cite{TW2014}. Define \(f_{00}^{(N)}\) as its onsite Fredholm
coefficient: the integral in \eqref{eq:double} with its global rational
factor replaced by \(1/(XY)\), retaining the same pair factors,
factorial and measures. The exact numerator identity
\[
 \frac{(1+X^{-1})(1+Y^{-1})}{(1-X)(1-Y)}
 =\frac1{XY}+\frac{2(X^{-1}+Y^{-1})}{(1-X)(1-Y)}
\]
therefore gives, without summing over particle number,
\[
 C_N^{\rm std}=f_{00}^{(N)}+2T_N.
\]
The local convex-gradient argument without either global product
factor bounds every fixed derivative of the onsite term at the
selected point. Smoothness is sufficient here; it is not silently
promoted to an analytic-extension theorem.

Boukraa et al.'s equations (5)--(9) \cite{Boukraa2008} give
\[
 C_N^{\rm std}=\frac1{N!}
 \int\prod_{i<N}\frac{\dd\varphi_i}{2\pi}
 \prod_i\frac1{\sinh\gamma_i}\,
 \frac{1+Z}{1-Z}
 \prod_{i<j}\left[
  \frac{\sin((\varphi_i-\varphi_j)/2)}
       {\sinh((\gamma_i+\gamma_j)/2)}
             \right]^2 ,
\]
where \(\sum_i\varphi_i=0\pmod{2\pi}\),
\(\cosh\gamma_i=S-\cos\varphi_i\), and \(z_i=e^{-\gamma_i}\).
The source variable \(w=s/[2(1+s^2)]\) satisfies \(1/(2w)=S\).
Its \(x_i\) is our \(z_i\), and its \(y_i\) is
\(1/\sinh\gamma_i\), not our contour coordinate \(y_i\).
These substitutions verify the observable, temperature side,
measures, and squared pair factor.

At the first selected order there are four signed ordered compact
charts, centered at \(+\alpha,-\alpha,+\beta,-\beta\).
Simultaneous angle reversal preserves both the density and real
integration measure. The \(\alpha/\beta\) exchange gives the same
coefficient by the sine-power calculation.
For any one chart the pair coefficient is
\(-1/(4\sin^2\beta)\), identical to its reduced lower-contour
counterpart. The product \(y_0^Nz_0^N=1\) removes the apparent
one-body phase discrepancy. The missing angular measure is
exactly the positive \(2\pi\) from the mean residue, so a single
standard chart has coefficient \(L_\beta\). Four standard charts
and two lower offsite charts agree with
\(C_N^{\rm std}=f_{00}^{(N)}+2T_N\).

\paragraph{Angle and power conventions in the review comparison.}
The printed amplitude in equation (69) of \cite{McCoyMaillard2012}
is not used to establish the internal coefficient, sign, or
multiplicity above. To make the comparison precise, put
\(a_0=(N^2-1)/2\) and \(\nu=a_0-1=k-1/2\).
The single-chart coefficient corresponding to an original-function
term \(A_{\rm chart}\eps^\nu\), modulo regular Taylor terms, is
\[
 A_{\rm chart}=L_\beta e^{\ii k\theta_*}
                      \frac{\sqrt\pi}{\Gamma(a_0)}.
\]
This follows by taking \(k\) derivatives with
\(\partial_s=e^{-\ii\theta_*}\partial_\eps\) along the radius.
For \(0<\theta_*<\pi/2\), the positive real factors in the local
formula leave the phase
\[
 \frac{A_{\rm chart}}{|A_{\rm chart}|}
      =(-1)^k e^{\ii\pi a_0/2}.
\]
Here \((-\ii d)^{-a_0}\) has the continuation specified in
Lemma~\ref{lem:period}. In the review's formula, substituting the
same positive angle and using the principal power of
\(N\ii\sin\theta_*\) gives phase \(e^{\ii\pi\nu/2}\).
The geometric denominator is positive for our selected angles.
The ratio of these phases is
\[
 (-1)^k e^{\ii\pi(a_0-\nu)/2}=-\ii,
\]
since \(k\) is odd. Substituting the negative angle instead gives
phase \(e^{-\ii\pi\nu/2}\) and phase ratio one. At \(|s_*|=1\),
that negative angle is \(\arg(s_*^{-1})\).

The review defines its angle through a cosine relation in (70);
that relation alone does not choose between these two signs or
identify its power branch with the exterior germ used here.
Consequently the last computation is a comparison under stated
conventions, not evidence that the review explicitly makes that
inverse-angle choice. Its low-temperature parameter is
\(\eps_{\rm rev}=1-s_*/s=\eps/(1+\eps)\); this real first-order
change cannot itself remove the phase difference. A full comparison
must also specify the continuation side, power branch, and single-chart
versus total-chart multiplicity. It is not justified to claim that
multiplicity is the only possible difference. None of these external
conventions enters the internal identity \(L_\alpha=L_\beta\).
No conclusion about the degenerate four-particle endpoint normalization
is imported into the present distinct, nonzero-angle family.

% END SOURCE


% BEGIN SOURCE: appendices/G_audit_map.tex
\section{Audit map, coverage, and status}\label{app:auditmap}
This appendix records the structure of the proposed proof. It is not
evidence for any estimate in the preceding sections. A correct
dependency map and a complete list of source files do not certify
the mathematics.

\subsection{External inputs and internally claimed estimates}
The external physical inputs are the spontaneous magnetization
\cite{Yang1952}, the low-temperature Fredholm correlation formula
\cite{PalmerTracy1981} in the representation of Appendix 1 of
\cite{TW2014}, and the fixed-even-order non-Nickel smoothness
argument of \cite{TW2014}. Appendix~\ref{app:normalization}
tracks the measures and compares the independent full-site formula
in \cite{Boukraa2008}. The review \cite{McCoyMaillard2012}
supplies physical definitions and amplitude conventions; its printed
amplitude is not used as a premise for the coefficient proved here.

The historical and methodological references have a separate role:
\cite{WMTB1976,GuttmannEnting1996,Nickel1999,Nickel2000}
locate the expansion and singularity problem;
\cite{TW2013,TW2017Toeplitz} locate the existing diagonal and
Toeplitz natural-boundary results and the noncancellation strategy.
None is invoked as a theorem giving the bulk uniform infinite-tail
estimate \eqref{eq:roadmaptail}. The bibliography records these
roles; it is not a certificate that every related publication has
been checked for priority.

The remaining analytic statements are proposed internal derivations
from the displayed formulas. In particular, the pair compactness
estimates, all protected disks, the high-order Lie bounds, and
the complete sector sum require adversarial checking.
The paper asserts a candidate proof of the full bulk statement;
it does not replace that statement by a diagonal susceptibility,
a finite sum, or a conditional numerical result.

\subsection{Dependency order}
The first-term route is
\[
 \begin{gathered}
 \text{normalization and residues}
 \ \longrightarrow\ \text{first-order arithmetic}\\
 \longrightarrow\ \text{mean residue and shape period}
 \ \longrightarrow\ \text{two-chart noncancellation}.
 \end{gathered}
\]
The tail route starts from the exact retraction and exterior normal
convergence. The selected-contour estimate uses its own \(c/N\)
disk. The pair/current route is ordered as
\[
 \begin{gathered}
 \text{radial-center complete-pair bound}
       \quad\text{(Lemma~\ref{lem:contraction})}\\
 \downarrow\\
 \text{original disk (Lemma~\ref{lem:originaldisk})}
 \quad\text{or}\quad
 \text{protected current disk (Lemma~\ref{lem:protected})}\\
 \downarrow\\
 \text{Cauchy bounds in Propositions~\ref{prop:highKS}
       and \ref{prop:largeS}.}
 \end{gathered}
\]
The radial-center lemma uses neither Cauchy bound.
For intermediate orders, the fixed support partition, the regular
branch geometry, and the recurrence of Lemma~\ref{lem:jets} precede
the microcore, all-branch complement, mixed, and compact-right
estimates. Those estimates, and not the final natural-boundary
criterion, are the inputs to Theorem~\ref{thm:tail}.

\subsection{Every region in the tail}
In the table, intermediate means \(N_0+2\le N<D\sqrt H\);
high means \(D\sqrt H\le N<C_H(H+1)^2\).
All rows apply to every \(0\le j\le k\). The \(N\)- and
\(\lambda\)-splits are made after differentiation.
\small
\begin{longtable}{>{\raggedright\arraybackslash}p{.06\textwidth}>{\raggedright\arraybackslash}p{.24\textwidth}>{\raggedright\arraybackslash}p{.39\textwidth}>{\raggedright\arraybackslash}p{.20\textwidth}}
\toprule
Row & Region & Required support or mechanism & Proof location\\
\midrule
\endfirsthead
\toprule
Row & Region & Required support or mechanism & Proof location\\
\midrule
\endhead
W1 & Ultra-high original &
Every angle, \(N\ge C_H(H+1)^2\); weighted matched pairs and
original disk. & Proposition~\ref{prop:ultrahigh}\\
W2 & Final selected \(F\) &
Named \(a_q>0,p_q=1\); protected branches; compact matching
penalty; disk \(c/N\). & Proposition~\ref{prop:F}\\
W3 & High \(K\) and small-current orders &
Enlarged lower support, with one named upper transition
variable for the current; original \(c\eps\) disk. &
Lemma~\ref{lem:originaldisk}, Proposition~\ref{prop:highKS}\\
W4 & Intermediate all-branch core &
Every \(|u_i|\le b_N\); resultant \(Y\)-gap and one-phase
branch transport. & Proposition~\ref{prop:micro}\\
W5 & All-branch far complement &
Named \(|u_a|\ge b_N/2\), \(\rho\ge\rho_N\); rational
weights and extreme-pair coarea after differentiation. &
Proposition~\ref{prop:allB}\\
W6 & All-branch near equality &
Same anchor, \(\rho\le2\rho_N\ll b_N\); common sign and
positive full-collision exponent. & Proposition~\ref{prop:allB}\\
W7 & Intermediate \(K\), compact left &
Outer anchor away from branch and a compact left root;
fixed \(Z\)-gap and one simple phase. & Proposition~\ref{prop:mixed}\\
W8 & Intermediate \(K\), mixed &
Outer compact \(q\), true-branch \(j\), \(|b_j|\ge4|b_q|\);
two-phase field at \(\lambda=0\). & Proposition~\ref{prop:mixed}\\
W9 & Intermediate \(K\), all right &
No true branch or left label; strict real phase curvature
and compact rational atlas. & Proposition~\ref{prop:compactR}\\
W10 & Intermediate small current, left &
Named \(q\) on an upper plateau and a left modulus gap;
all moving branches remain on true plateaus. &
Proposition~\ref{prop:mixed}\\
W11 & Intermediate small current, mixed &
Upper-right \(q\), true-branch \(j\); exact preservation of
both occupancy sums in every derivative. &
Proposition~\ref{prop:mixed}\\
W12 & Intermediate small current, all right &
Only compact right variables; coupled Hessian perturbation
\(N^{C_j}(\eps+\lambda_*)\to0\). &
Proposition~\ref{prop:compactR}\\
W13 & Large current &
Named \(q\), \(P\ge1\), \(\lambda\ge\lambda_*\);
disk \(c\lambda_*/N^2\), root inflation and product gaps. &
Lemma~\ref{lem:protected}, Proposition~\ref{prop:largeS}\\
\bottomrule
\end{longtable}
\normalsize

The upper inverse branch \(+\theta_B\) is in the \(a=p=1\)
plateau. It can appear in \(F\), where it is shifted inward;
it cannot appear on a \(K\) or Stokes transition support.
The nested partition in Lemma~\ref{lem:partition} keeps every
derivative of a compact label out of the inner branch region.
These are support assertions required for coverage, not
optional choices after an estimate has been applied.

\subsection{Load-bearing checks for an adversarial reader}
The most direct way to disprove the proposed argument is to
find a failure of one of the two estimates
\eqref{eq:roadmapfirst}--\eqref{eq:roadmaptail}.
The following checks expose their sensitive interfaces:
\begin{enumerate}
\item Verify the physical measure count, the clockwise mean
closure and its positive \(2\pi\), and the complex phase
equality \(L_\alpha=L_\beta\). Check the complement of the
two charts locally; the selected point itself is Nickel.
\item Verify same-group strictness and the cross-group and named-pair
bounds on the actual supports, including endpoint limits, a moving
lower branch, the named upper variable, and each stated
complex disk. Do not infer this from separate Schur bounds
where one such factor has a pole.
\item Check \eqref{eq:pullbackLie} and every case in
\eqref{eq:jetdegree} at arbitrary fixed order, retaining the
unfactored numerator and the real cutoff derivatives.
Check both kinds of puncture flux before using integration
by parts repeatedly.
\item Check the selected branch measure divided by the real
coarea Jacobian, the multiplicity of each phase map, and
the exponential scales \(b_N,\rho_N\).
An uncontrolled power of \(\eps^{-1}\), an unabsorbed
\(e^{cN^3}\), or a loss of strict pair contraction would
invalidate the stated window sum.
\item Check that the thirteen rows cover every derivative
density and that their constants can be selected in the
order of Appendix~\ref{app:constants}. The argument needs
all \(j\le k\), with \(k\) fixed by the selected point.
\end{enumerate}

\subsection{Reconstruction boundary}
Revision 2 retains the first nonzero-angle prime family and
the final thirteen-sector tail route of checkpoint v38.
It expands the real-cutoff jet recurrence, separates the
radial-center and protected-disk arguments, and exposes the
hybrid measure and constant choices. These additions are
mathematical exposition to audit, not an external validation.

Revision 3 corrects the exterior convergence paragraph in
Appendix~\ref{app:pfaffian}: global inverse products are bounded
exponentially in \(N\), rather than independently of \(N\).
The displayed Pfaffian--factorial bound absorbs this cost, and
the compact-set Cauchy estimate makes the differentiated
normal convergence explicit. This revision leaves the main
claim, boundary-point family, and tail estimates unchanged.

Revision 4 corrects the one-body argument in
Lemma~\ref{lem:originaldisk}: the branch-center displacement
on a complex parameter disk includes \(|s-\se|\).
The two-region estimate now retains this perturbation and
proves the same occupancy-independent integrable majorant.
No contour, particle-number window, or main claim is changed.

Revision 5 incorporates the explicitly identified first-term gluing
supplement from the subsequent full-chain audit. Section~\ref{sec:first}
now constructs a \(y\)-only angular partition before taking any
\(x\) residues, replaces its smooth mean cutoffs only across separated
edge regions, and gives a quantitative finite-cover estimate for the
complement. These additions connect the local residue calculation to
the whole double torus. They are new exposition and supporting
derivations relative to revision 4, not a claim that revision 4 already
contained this construction. The tail estimates and the main target
are unchanged; the audit's other verification labels are not imported
as premises.

Revision 6 makes four further interfaces explicit following a
separate audit of revision 5: compact upper supports and branch
plateaus; a periodic left/right partition whose extra seam lies
where all relevant weighted jets vanish; Cauchy's inequality
applied to the complete holomorphic \(F_N\); and exponential-scale
inclusions imposed at every finite tail order. The related negative
anchor scale is chosen with the same uniform inclusion. These are
explicit choices and local justifications within the existing proof
route. The model, theorem, point family, derivative range, and
particle-number windows are unchanged. No audit verdict is used
as a mathematical premise.

Revision 7 corrects the logarithmic expansion of the microcore
majorant by displaying its negative \(N^2\log N/2\) term. It also
records the angle and power conventions required for a comparison
with the review amplitude, and explicitly includes the unwrapped
phase lengths and period counts in near-equality coarea estimates.
The physical-state convention is stated at the outset; the current
E1 integration specifies the \(+\) boundary-condition state. These are local corrections and
clarifications; the two main estimates, point family, derivative
range, and particle windows are unchanged. The received audit's
failure to find a counterexample is not an additional proof premise.

Revision 8 supplies the missing historical and methodological
attribution, including the diagonal noncancellation argument and the
upstream correlation representation. It distinguishes the proposed
bulk estimates from the existing proof architecture. The endpoint
wording in Lemma~\ref{lem:contraction} now asserts strictness only
within groups and explicitly permits cross-group equality. Inverse
temperature and the inner branch cutoff have distinct symbols from
the geometric angle \(\beta\). The selected family, two central
estimates, derivative range, and particle windows are unchanged.
The separate revision record states the checks actually performed;
earlier audit verdicts are attached to their earlier manuscript
versions and are not automatically new audits of revision 8.

Earlier routes involving a jointly absolute mean/shape
limit, a multiplier obtained by dividing by the original
Vandermonde, a central connector contour, or a degenerate
four-particle endpoint amplitude are not assumptions of
this manuscript. The separate source-coverage record names
their replacements in v38. A historical file being
present in the checkpoint is not by itself a reason to
include its superseded argument in the proof.

The accompanying release records identify file hashes, the active-source
mapping, and the edits from the preceding reconstruction.
Those checks establish provenance and textual coverage.
They do not establish truth of any proposition, settle
literature priority, or constitute a completed independent
adversarial audit.

% END SOURCE

\subsection*{Rc4 W2 repair}
Version 0.1-rc4 adds the explicit support geometry,
local root-continuation and complete-pair transfer arguments,
reserved disk slack, tau-dependent protected-disk restrictions,
and coupled-occupancy product-rule count. Lemma~\ref{lem:originaldisk}
now explicitly confines its pair estimates to \(K\) and named
Stokes supports. The theorem, point family, derivative range and
particle windows are unchanged. These are new scoped proof repairs,
not arguments attributed retrospectively to rc3. The supplied
audits and finite diagnostics were treated as untrusted evidence.
This candidate is prepared for public review. The subsequent E1 integration
checks the primary-source interface, specifies the \(+\) state, and
reconciles the published full-site coefficient with the existing
offsite formulas. It changes no W1/W2 estimate or theorem statement.
No human, peer, or proof-assistant certification is claimed.

\subsection*{Release status and AI assistance}
Version 0.1-rc3 incorporates the W1 proof-interface repairs described in this revision record.
It adds the current support partition and deformed branch geometry,
recounts the remaining pair factors after differentiation, corrects
the near-collision slope cost, and expands the constant dependencies.
These substantive proof repairs received a scoped AI analytic audit;
the historical rc2 records retain their original inputs and verdicts.
Three problem-only reconstruction runs returned partial results;
none reconstructed the complete natural-boundary proof. One
method-informed run claimed the full tail estimate while explicitly
leaving its whole-integral first singular term (node I3) unresolved.
Its tail claim has not received a complete independent verification
by the receiving workflow. The different point families used in
some runs do not allow their conclusions to be combined silently.
The repository separates these run claims from manuscript claims,
version-specific audit reports, and finite computational checks.

Generative AI was used extensively in research exploration,
mathematical derivation, manuscript writing, code, internal audits,
and reconstruction attempts. Dehao Lin is the named human author
and public contact. No specific division of mathematical contributions
or line-by-line human verification is asserted here.
An AI audit verdict is not proof certification. No independent human
expert certification or proof-assistant verification is claimed.
The author declares no additional funding or competing interests
for this project. Project text is offered under CC BY 4.0 to the extent
the author holds rights in it; project code is offered under the MIT
license. Cited third-party works retain their own rights.
The accompanying provenance record states
the limits of the available historical access evidence.

\begin{thebibliography}{99}

\bibitem{Boukraa2008}
S.~Boukraa, A.~J. Guttmann, S.~Hassani, I.~Jensen, J.-M. Maillard,
B.~Nickel, and N.~Zenine.
\newblock Experimental mathematics on the magnetic susceptibility of the
square lattice {Ising} model.
\newblock \emph{Journal of Physics A: Mathematical and Theoretical},
41(45):455202, 2008.
\newblock Formula comparison uses equations (4)--(9) of
\href{https://arxiv.org/abs/0808.0763v1}{arXiv:0808.0763v1}.

\bibitem{GuttmannEnting1996}
A.~J. Guttmann and I.~G. Enting.
\newblock Solvability of some statistical mechanical systems.
\newblock \emph{Physical Review Letters}, 76:344--347, 1996.
\newblock \href{https://doi.org/10.1103/PhysRevLett.76.344}{doi:10.1103/PhysRevLett.76.344}.

\bibitem{McCoyMaillard2012}
Barry~M. McCoy and Jean-Marie Maillard.
\newblock The importance of the {Ising} model.
\newblock \emph{Progress of Theoretical Physics}, 127(5):791--817, 2012.
\newblock \href{https://arxiv.org/abs/1203.1456}{arXiv:1203.1456}.
Physical definitions and amplitude conventions; the printed amplitude
is not used as a premise for the internal coefficient.

\bibitem{Nickel1999}
B.~Nickel.
\newblock On the singularity structure of the 2D {Ising} model susceptibility.
\newblock \emph{Journal of Physics A: Mathematical and General},
32(21):3889--3906, 1999.
\newblock \href{https://doi.org/10.1088/0305-4470/32/21/303}{doi:10.1088/0305-4470/32/21/303}.

\bibitem{Nickel2000}
B.~Nickel.
\newblock Addendum to ``On the singularity structure of the 2D {Ising}
model susceptibility''.
\newblock \emph{Journal of Physics A: Mathematical and General},
33(8):1693--1711, 2000.
\newblock \href{https://doi.org/10.1088/0305-4470/33/8/313}{doi:10.1088/0305-4470/33/8/313}.

\bibitem{Orrick2001}
W.~P. Orrick, B.~Nickel, A.~J. Guttmann, and J.~H.~H. Perk.
\newblock The susceptibility of the square lattice {Ising} model:
New developments.
\newblock \emph{Journal of Statistical Physics}, 102:795--841, 2001.
\newblock \href{https://doi.org/10.1023/A:1004850919647}{doi:10.1023/A:1004850919647};
\href{https://arxiv.org/abs/cond-mat/0103074v1}{arXiv:cond-mat/0103074v1}.

\bibitem{PalmerTracy1981}
John Palmer and Craig Tracy.
\newblock Two-dimensional {Ising} correlations: Convergence of the scaling limit.
\newblock \emph{Advances in Applied Mathematics}, 2(3):329--388, 1981.
\newblock \href{https://doi.org/10.1016/0196-8858(81)90010-5}{doi:10.1016/0196-8858(81)90010-5}.

\bibitem{TW2013}
Craig~A. Tracy and Harold Widom.
\newblock On the diagonal susceptibility of the two-dimensional {Ising} model.
\newblock \emph{Journal of Mathematical Physics}, 54:123302, 2013.
\newblock \href{https://arxiv.org/abs/1307.5913v3}{arXiv:1307.5913v3}.

\bibitem{TW2014}
Craig~A. Tracy and Harold Widom.
\newblock On the singularities in the susceptibility expansion for the
two-dimensional {Ising} model.
\newblock \emph{Journal of Statistical Physics}, 156(6):1125--1135, 2014.
\newblock \href{https://doi.org/10.1007/s10955-014-1061-4}{doi:10.1007/s10955-014-1061-4}.
\newblock Formulas are cited from this published version.

\bibitem{TW2017Toeplitz}
Craig~A. Tracy and Harold Widom.
\newblock Natural boundary for a sum involving {Toeplitz} determinants.
\newblock In \emph{Large Truncated Toeplitz Matrices, Toeplitz Operators,
and Related Topics}, pages 703--718. Springer, 2017.
\newblock Preprint (2015):
\href{https://arxiv.org/abs/1502.04922}{arXiv:1502.04922}.

\bibitem{WMTB1976}
Tai Tsun Wu, Barry~M. McCoy, Craig~A. Tracy, and Eytan Barouch.
\newblock Spin-spin correlation functions for the two-dimensional {Ising}
model: Exact theory in the scaling region.
\newblock \emph{Physical Review B}, 13:316--374, 1976.
\newblock \href{https://doi.org/10.1103/PhysRevB.13.316}{doi:10.1103/PhysRevB.13.316}.

\bibitem{Yang1952}
C.~N. Yang.
\newblock The spontaneous magnetization of a two-dimensional {Ising} model.
\newblock \emph{Physical Review}, 85:808--816, 1952.
\newblock \href{https://doi.org/10.1103/PhysRev.85.808}{doi:10.1103/PhysRev.85.808}.

\end{thebibliography}

\end{document}
